% =============================================================================
% MenGrowth-Model: Foundation Model Pipeline for Meningioma Growth Forecasting
% BSc Thesis — Technical Report
% =============================================================================
\documentclass[12pt,a4paper]{article}

% --- Encoding and fonts ---
\usepackage[utf8]{inputenc}
\usepackage[T1]{fontenc}
\usepackage{lmodern}

% --- Mathematics ---
\usepackage{amsmath}
\usepackage{amssymb}
\usepackage{amsthm}
\usepackage{mathtools}
\usepackage{bm}

% --- Tables and figures ---
\usepackage{graphicx}
\usepackage{booktabs}
\usepackage{multirow}
\usepackage{subcaption}
\usepackage{float}

% --- Layout ---
\usepackage[margin=2.5cm]{geometry}
\usepackage{setspace}
\onehalfspacing

% --- References and links ---
\usepackage[colorlinks=true, linkcolor=blue!70!black, citecolor=green!50!black, urlcolor=blue!60!black]{hyperref}
\usepackage[capitalise,noabbrev]{cleveref}

% --- Code listings ---
\usepackage{listings}
\lstset{
  basicstyle=\ttfamily\small,
  breaklines=true,
  frame=single,
  numbers=left,
  numberstyle=\tiny\color{gray},
}

% --- Algorithms ---
\usepackage{algorithm}
\usepackage{algorithmic}

% --- Misc ---
\usepackage{xcolor}
\usepackage{enumitem}
\usepackage{siunitx}

% --- Theorem environments ---
\newtheorem{definition}{Definition}
\newtheorem{proposition}{Proposition}
\newtheorem{remark}{Remark}

% --- Custom commands ---
\newcommand{\R}{\mathbb{R}}
\newcommand{\E}{\mathbb{E}}
\newcommand{\Vol}{\mathbf{z}_{\text{vol}}}
\newcommand{\Loc}{\mathbf{z}_{\text{loc}}}
\newcommand{\Shape}{\mathbf{z}_{\text{shape}}}
\newcommand{\Res}{\mathbf{z}_{\text{res}}}
\newcommand{\dCor}{\operatorname{dCor}}
\newcommand{\dCov}{\operatorname{dCov}}
\newcommand{\dVar}{\operatorname{dVar}}
\newcommand{\KL}{\operatorname{KL}}
\newcommand{\diag}{\operatorname{diag}}
\newcommand{\tr}{\operatorname{tr}}
\newcommand{\LOPO}{\text{LOPO}}
\newcommand{\MOGP}{\text{MOGP}}
\newcommand{\PAMOGP}{\text{PA-MOGP}}
\newcommand{\LME}{\text{LME}}
\newcommand{\HGP}{\text{H-GP}}
\newcommand{\GP}{\mathcal{GP}}
\newcommand{\LoRA}{\text{LoRA}}
\newcommand{\SDP}{\text{SDP}}

% =============================================================================
% Title
% =============================================================================
\title{%
  \textbf{Foundation Model Adaptation for Meningioma Growth Forecasting\\
  from Multi-Modal 3D MRI}\\[0.8em]
  \large A Gaussian Process Hierarchy over Disentangled Latent Trajectories
}

\author{%
  Mario Pascual-Gonz\'{a}lez\\
  \textit{Department of Computer Science}\\
  \textit{Universidad de Granada}\\
  \texttt{mariopascual@ugr.es}
}

\date{March 2026}

% =============================================================================
% Document
% =============================================================================
\begin{document}

\maketitle

% --- Abstract ---
\begin{abstract}
Meningiomas are the most prevalent primary intracranial tumours, yet predicting
their growth trajectories from longitudinal MRI remains an open challenge due to
heterogeneous growth dynamics and limited longitudinal data. We present a
four-phase pipeline that adapts BrainSegFounder---a SwinUNETR foundation model
pretrained on 41,400+ brain MRI subjects---to the meningioma domain. Phase~1
applies Low-Rank Adaptation (LoRA) to the encoder's Stages~3--4 using
segmentation as a proxy task on BraTS-MEN~2024 (1,000~subjects). Phase~2 trains
a Supervised Disentangled Projection (SDP) that maps frozen encoder features
$\mathbf{h} \in \R^{768}$ to a structured latent space
$\mathbf{z} \in \R^{128}$ with semantically supervised partitions for tumour
volume, location, and shape. Phase~3 encodes a longitudinal Andalusian cohort
(42~patients, 137~studies) through the frozen pipeline with optional ComBat
harmonisation. Phase~4 fits a three-model Gaussian process hierarchy---Linear
Mixed-Effects (LME), Hierarchical GP (H-GP), and Partition-Aware Multi-Output
GP (PA-MOGP)---evaluated via Leave-One-Patient-Out Cross-Validation across
33~patients. The pipeline replaces a previously considered Neural ODE approach,
which was abandoned due to catastrophic overparameterisation (3,100+~parameters
for 112~observation pairs). We report domain gap analysis, LoRA ablation
results including dual-domain adaptation experiments, and discuss methodological
revisions motivated by empirical findings including the failure of shape
disentanglement and the temporal stationarity of meningioma location.
\end{abstract}

\newpage
\tableofcontents
\newpage

% --- Section 1: Introduction ---
% =============================================================================
% Section 1: Introduction
% =============================================================================
\section{Introduction}
\label{sec:introduction}

% ---------------------------------------------------------------------------
% 1.1 Clinical Motivation
% ---------------------------------------------------------------------------
\subsection{Clinical Motivation}
\label{sec:intro:clinical}

Meningiomas are the most common primary intracranial tumours, accounting for
approximately 39\% of all central nervous system neoplasms, with an
age-adjusted incidence rate of 9.51 per 100,000 population
\cite{ostrom2023brain}. The majority are classified as WHO Grade~I (benign),
and a substantial fraction are discovered incidentally on neuroimaging
performed for unrelated indications \cite{vernooij2007incidental}. For these
incidentally discovered meningiomas, clinical management hinges on a
deceptively simple question: \emph{will this tumour grow?}

Current clinical practice relies on serial MRI surveillance with manual
volumetric assessment at intervals of 3--12 months
\cite{hashimoto2012natural,ius2020management}. This approach has three
fundamental limitations. First, manual volumetry is labour-intensive and
subject to considerable inter-observer variability, with reported coefficient
of variation exceeding 15\% for small tumours \cite{defined2016interobserver}.
Second, the decision to intervene (surgery or radiotherapy) is typically made
reactively---after significant growth has already occurred---rather than
proactively based on predicted trajectories. Third, meningioma growth dynamics
are heterogeneous: some tumours remain stable for decades, others exhibit
steady linear growth, and a minority display accelerating or saltatory growth
patterns \cite{hashimoto2012natural,nakamura2003natural}. No established
imaging biomarker reliably predicts which pattern a given tumour will follow.

Quantitative growth prediction from longitudinal MRI could transform
meningioma management by enabling personalised surveillance schedules,
early identification of aggressive growth patterns, and evidence-based
timing of intervention. However, building such predictive models faces a
data bottleneck: longitudinal meningioma MRI datasets are rare, typically
small (tens of patients), and collected retrospectively across heterogeneous
scanners and protocols.

% ---------------------------------------------------------------------------
% 1.2 Technical Challenge
% ---------------------------------------------------------------------------
\subsection{Technical Challenge}
\label{sec:intro:technical}

The central technical challenge is to build a growth prediction model from a
longitudinal cohort of only 33~patients with 2--6~MRI studies each (100~total
observations), while leveraging cross-sectional data from 1,000+ subjects to
learn rich tumour representations. This requires solving three interrelated
problems:

\begin{enumerate}[leftmargin=*]
  \item \textbf{Domain adaptation.} Existing brain MRI foundation models are
    pretrained on glioma data, which differs fundamentally from meningioma in
    tumour morphology, location, and tissue composition. Direct application
    of a glioma-trained encoder to meningioma data produces suboptimal feature
    representations (see \cref{sec:domain_gap}).

  \item \textbf{Structured representation learning.} Growth prediction
    requires features that are not only discriminative but also
    \emph{semantically structured}: tumour volume, spatial location, and
    morphological shape should be disentangled into separate latent
    subspaces amenable to interpretable temporal modelling.

  \item \textbf{Small-sample temporal modelling.} With 33~patients and sparse,
    irregularly-sampled time series, conventional deep learning approaches
    (recurrent networks, Neural ODEs) are catastrophically overparameterised.
    The modelling framework must provide calibrated uncertainty and degrade
    gracefully with limited data.
\end{enumerate}

% ---------------------------------------------------------------------------
% 1.3 Proposed Approach
% ---------------------------------------------------------------------------
\subsection{Proposed Approach}
\label{sec:intro:approach}

We address these challenges through a four-phase pipeline that progressively
transforms raw multi-modal 3D MRI into patient-specific growth predictions
with calibrated uncertainty:

\begin{description}[leftmargin=*,labelwidth=!]
  \item[Phase 1 --- LoRA Adaptation (\cref{sec:lora}).] We adapt
    BrainSegFounder \cite{cox2024brainsegfounder}, a SwinUNETR
    \cite{hatamizadeh2022swinunetr} foundation model pretrained via
    self-supervised learning on 41,400+ brain MRI subjects, to the meningioma
    domain using Low-Rank Adaptation (LoRA) \cite{hu2022lora}. LoRA injects
    trainable low-rank matrices into the encoder's Stages~3--4, adapting
    high-level features while preserving low-level anatomical representations.
    Segmentation on BraTS-MEN~2024 \cite{labella2024brats} serves as the
    proxy task, with optional auxiliary semantic heads providing multi-task
    learning signal.

  \item[Phase 2 --- Supervised Disentangled Projection (\cref{sec:sdp}).] A
    lightweight projection network maps frozen encoder features
    $\mathbf{h} \in \R^{768}$ to a structured latent space
    $\mathbf{z} \in \R^{128}$ partitioned into semantically supervised
    subspaces for volume, location, shape, and a residual component.
    Training combines semantic regression, VICReg-style
    \cite{bardes2022vicreg} variance--covariance regularisation, and distance
    correlation \cite{szekely2007measuring} penalties via a four-phase
    curriculum.

  \item[Phase 3 --- Cohort Encoding (\cref{sec:encoding}).] All longitudinal
    MRI studies from the Andalusian cohort are encoded through the frozen
    pipeline, producing per-patient latent trajectories. ComBat harmonisation
    \cite{johnson2007adjusting,fortin2018harmonization} is assessed and
    applied if significant scanner-induced distributional shift is detected.

  \item[Phase 4 --- Growth Prediction (\cref{sec:growth}).] A three-model
    Gaussian process hierarchy of increasing complexity---Linear Mixed-Effects
    (LME), Hierarchical GP (H-GP), and Partition-Aware Multi-Output GP
    (PA-MOGP)---is fitted on the latent trajectories and evaluated via
    Leave-One-Patient-Out Cross-Validation (LOPO-CV) across 33~folds.
    Volume predictions are decoded through the frozen SDP semantic head and
    compared against observed tumour volumes.
\end{description}

% ---------------------------------------------------------------------------
% 1.4 Contributions
% ---------------------------------------------------------------------------
\subsection{Contributions}
\label{sec:intro:contributions}

This thesis makes the following contributions:

\begin{enumerate}[leftmargin=*]
  \item \textbf{Domain adaptation analysis.} We provide a quantitative
    characterisation of the glioma--meningioma domain gap in the
    BrainSegFounder feature space using Maximum Mean Discrepancy, Centred
    Kernel Alignment, and semantic linear probes (\cref{sec:domain_gap}).

  \item \textbf{LoRA adaptation with semantic supervision.} We demonstrate
    LoRA adaptation of a 3D medical imaging foundation model for meningioma
    segmentation, with ablations across rank, adapter type (LoRA vs.\ DoRA),
    auxiliary semantic heads, and single- vs.\ dual-domain training
    (\cref{sec:lora}).

  \item \textbf{Supervised Disentangled Projection.} We design and
    theoretically justify a non-generative approach to structured latent space
    construction, motivated by the impossibility of unsupervised
    disentanglement \cite{locatello2019challenging} and the failure of
    VAE-based alternatives in our experimental setting (\cref{sec:sdp}).

  \item \textbf{GP hierarchy for latent growth prediction.} We propose a
    principled replacement for Neural ODE-based temporal modelling, with
    closed-form inference, calibrated uncertainty, and a parameter count
    appropriate for the data regime (6--95 parameters vs.\ 3,100+ for the
    Neural ODE) (\cref{sec:growth}).

  \item \textbf{Empirical methodology revision.} We document how empirical
    findings---including shape disentanglement failure ($R^2 \leq 0.11$),
    temporal stationarity of meningioma location, and marginal LoRA benefit
    ($\Delta R^2 = +0.016$)---motivated principled revisions to the pipeline
    architecture (\cref{sec:evaluation}).
\end{enumerate}

% ---------------------------------------------------------------------------
% 1.5 Document Structure
% ---------------------------------------------------------------------------
\subsection{Document Structure}

The remainder of this report is organised as follows.
\Cref{sec:background} reviews the theoretical foundations underlying each
pipeline component. \Cref{sec:data} describes the datasets, preprocessing,
and data splits. \Cref{sec:domain_gap} presents the domain gap analysis
motivating LoRA adaptation. \Cref{sec:lora,sec:sdp,sec:encoding,sec:growth}
detail Phases~1--4 respectively. \Cref{sec:evaluation} consolidates quality
targets, ablation results, and model comparison. \Cref{sec:conclusion}
discusses limitations and future directions.


% --- Section 2: Background and Related Work ---
% =============================================================================
% Section 2: Background and Related Work
% =============================================================================
\section{Background and Related Work}
\label{sec:background}

% ---------------------------------------------------------------------------
% 2.1 Foundation Models for Brain MRI
% ---------------------------------------------------------------------------
\subsection{Foundation Models for Brain MRI}
\label{sec:bg:foundation}

The paradigm of pretraining large encoder--decoder architectures on massive,
diverse datasets and subsequently fine-tuning on downstream tasks has
transformed medical image analysis \cite{he2023foundation,moor2023foundation}.
BrainSegFounder \cite{cox2024brainsegfounder} instantiates this paradigm for
3D brain MRI segmentation through a three-stage training pipeline:

\begin{enumerate}[leftmargin=*]
  \item \textbf{Self-supervised pretraining.} A SwinUNETR encoder
    \cite{hatamizadeh2022swinunetr,tang2022self} is pretrained on 41,400+
    brain MRI subjects from the UK Biobank using a combination of masked image
    modelling, contrastive learning, and rotation prediction. This stage uses a
    $96^3$ voxel region-of-interest (ROI) and produces general-purpose
    anatomical features.

  \item \textbf{Supervised pretraining.} The encoder is paired with the
    SwinUNETR decoder and trained on the BraTS~2021 glioma segmentation
    benchmark \cite{baid2021rsna} at $128^3$ ROI, producing tumour-specific
    features.

  \item \textbf{Per-fold fine-tuning.} Final fine-tuning on BraTS~2021 with
    cross-validation yields the released checkpoints.
\end{enumerate}

\subsubsection{SwinUNETR Architecture}

SwinUNETR combines the Swin Transformer \cite{liu2021swin} with a
UNet-style encoder--decoder structure. The encoder consists of a patch
embedding layer followed by four transformer stages with shifted window
self-attention. For BSF-Tiny (the variant used in this work), the
architectural parameters are:

\begin{table}[H]
  \centering
  \caption{SwinUNETR encoder (BSF-Tiny) architecture. Input: $[B, 4, 128^3]$.}
  \label{tab:swinunetr_arch}
  \begin{tabular}{lcccc}
    \toprule
    \textbf{Stage} & \textbf{Output shape} & \textbf{Channels} & \textbf{Heads} & \textbf{Blocks} \\
    \midrule
    Patch embed & $[B, 48, 64^3]$  & 48  & --- & --- \\
    Stage 1     & $[B, 96, 32^3]$  & 96  & 3   & 2 \\
    Stage 2     & $[B, 192, 16^3]$ & 192 & 6   & 2 \\
    Stage 3     & $[B, 384, 8^3]$  & 384 & 12  & 2 \\
    Stage 4     & $[B, 768, 4^3]$  & 768 & 24  & 2 \\
    encoder10   & $[B, 768, 4^3]$  & 768 & --- & --- \\
    \bottomrule
  \end{tabular}
\end{table}

The bottleneck layer \texttt{encoder10} applies a residual convolutional
block preserving the $768$-channel, $4^3$-spatial resolution from Stage~4.
Global Average Pooling (GAP) over the spatial dimensions produces the feature
vector $\mathbf{h} \in \R^{768}$ used throughout our pipeline.

The segmentation decoder mirrors the encoder with skip connections, producing
a 3-channel sigmoid output following the BraTS convention: Channel~0 encodes
Tumour Core (TC), Channel~1 encodes Whole Tumour (WT), and Channel~2 encodes
Enhancing Tumour (ET), using hierarchical overlapping labels rather than
mutually exclusive classes.

% ---------------------------------------------------------------------------
% 2.2 Low-Rank Adaptation (LoRA)
% ---------------------------------------------------------------------------
\subsection{Low-Rank Adaptation}
\label{sec:bg:lora}

Low-Rank Adaptation (LoRA) \cite{hu2022lora} provides parameter-efficient
fine-tuning by injecting trainable rank-$r$ decompositions into frozen
pretrained weight matrices. For a pretrained weight matrix
$\mathbf{W}_0 \in \R^{d \times k}$, LoRA parameterises the adapted weight as:
\begin{equation}
  \mathbf{W} = \mathbf{W}_0 + \frac{\alpha}{r}\, \mathbf{B}\mathbf{A},
  \quad \mathbf{B} \in \R^{d \times r},\; \mathbf{A} \in \R^{r \times k},
  \label{eq:lora}
\end{equation}
where $r \ll \min(d, k)$ is the rank, $\alpha$ is a scaling hyperparameter,
$\mathbf{A}$ is initialised from $\mathcal{N}(0, \sigma^2)$, and
$\mathbf{B} = \mathbf{0}$ so that the adapted model initially reproduces the
pretrained model exactly. The number of trainable parameters per adapted layer
is $r(d + k)$, compared to $dk$ for full fine-tuning.

The theoretical foundation rests on the observation by Aghajanyan et al.\
\cite{aghajanyan2021intrinsic} that pretrained language models have a low
intrinsic dimensionality: the parameter space can be projected to a
low-dimensional subspace without significant loss in fine-tuning performance.
LoRA operationalises this insight by constraining the weight update to lie in a
rank-$r$ subspace.

\textbf{DoRA} \cite{liu2024dora} extends LoRA by decomposing the weight
update into magnitude and direction components:
\begin{equation}
  \mathbf{W} = m \cdot \frac{\mathbf{W}_0 + \mathbf{B}\mathbf{A}}
    {\|\mathbf{W}_0 + \mathbf{B}\mathbf{A}\|_c},
  \label{eq:dora}
\end{equation}
where $m$ is a learnable magnitude vector and $\|\cdot\|_c$ denotes
column-wise normalisation. This decomposition was shown to more closely
approximate full fine-tuning behaviour in language models; we evaluate it as
an ablation condition (A2, \cref{sec:eval:ablations}).

% ---------------------------------------------------------------------------
% 2.3 Supervised Disentangled Projection
% ---------------------------------------------------------------------------
\subsection{Disentangled Representation Learning}
\label{sec:bg:disentanglement}

Disentangled representations encode independent generative factors of
variation in separate latent dimensions or subspaces. Higgins et al.\
\cite{higgins2018towards} formalised this via group theory: a representation
is disentangled with respect to a group decomposition
$G = G_1 \times \cdots \times G_n$ if there exists a corresponding
decomposition $Z = Z_1 \oplus \cdots \oplus Z_n$ where each $Z_i$ is affected
only by $G_i$.

The impossibility theorem of Locatello et al.\ \cite{locatello2019challenging}
established that unsupervised disentanglement is theoretically impossible
without inductive biases: for any generative model producing disentangled
representations, there exists another model with identical marginal
distribution but entangled latents. Their follow-up work
\cite{locatello2020disentangling} demonstrated that as few as 100 labelled
examples suffice to resolve this impossibility through supervision.

These results motivate our Supervised Disentangled Projection (SDP) approach:
rather than relying on VAE-based unsupervised disentanglement (which we
empirically found to suffer from posterior collapse and KL distortion), we
use explicit semantic supervision from segmentation-derived features to
anchor each latent partition to a known generative factor.

\subsubsection{VICReg Regularisation}

VICReg \cite{bardes2022vicreg} defines three regularisation terms for
representation learning. The \textit{variance} term prevents dimensional
collapse via a hinge loss on per-dimension standard deviation:
\begin{equation}
  \mathcal{L}_{\text{var}}(\mathbf{Z}) = \frac{1}{d}\sum_{j=1}^{d}
    \max\!\bigl(0,\; \gamma - \sqrt{\operatorname{Var}(z^{(j)}) + \epsilon}
    \,\bigr),
  \label{eq:vicreg_var}
\end{equation}
where $\gamma = 1$ is the target standard deviation. The \textit{covariance}
term decorrelates dimensions by penalising off-diagonal elements of the batch
covariance matrix $\mathbf{C}$:
\begin{equation}
  \mathcal{L}_{\text{cov}}(\mathbf{Z}) = \frac{1}{d}\sum_{i \neq j}
    C_{ij}^2.
  \label{eq:vicreg_cov}
\end{equation}

\subsubsection{Distance Correlation}

VICReg's covariance penalty captures only linear dependencies. Distance
correlation \cite{szekely2007measuring} detects all forms of statistical
dependence between random vectors $X$ and $Y$ of potentially different
dimensions:
\begin{equation}
  \dCor(X, Y) = \frac{\dCov(X, Y)}
    {\sqrt{\dVar(X) \cdot \dVar(Y)}},
  \label{eq:dcor}
\end{equation}
where $\dCov$ is computed from doubly-centred pairwise distance matrices.
Critically, $\dCor(X, Y) = 0$ if and only if $X \perp\!\!\!\perp Y$ (for
random vectors with finite first moments), providing a necessary and
sufficient condition for statistical independence.

\subsubsection{Spectral Normalisation}

Spectral normalisation \cite{miyato2018spectral} constrains each linear layer
$\mathbf{W}$ by dividing by its largest singular value
$\sigma(\mathbf{W})$, yielding a weight matrix with spectral norm exactly~1.
For an $L$-layer MLP with 1-Lipschitz activations (e.g., ReLU, GELU), this
bounds the global Lipschitz constant: $\operatorname{Lip}(f) \leq
\prod_{\ell=1}^{L} 1 = 1$. The contraction property
$\|\mathbf{z}_1 - \mathbf{z}_2\| \leq \|\mathbf{h}_1 - \mathbf{h}_2\|$
ensures that encoder feature noise is not amplified through the projection,
which is essential for stable downstream GP prediction
\cite{gouk2021regularisation}.

% ---------------------------------------------------------------------------
% 2.4 Gaussian Process Regression
% ---------------------------------------------------------------------------
\subsection{Gaussian Process Regression}
\label{sec:bg:gp}

A Gaussian process (GP) is a distribution over functions such that any finite
collection of function values is jointly Gaussian
\cite{rasmussen2006gaussian}:
\begin{equation}
  f(\mathbf{t}) \sim \GP\bigl(m(\mathbf{t}),\; k(\mathbf{t}, \mathbf{t}')\bigr),
  \label{eq:gp_def}
\end{equation}
where $m(\mathbf{t})$ is the mean function and $k(\mathbf{t}, \mathbf{t}')$
is the covariance (kernel) function. Given $n$ observations
$\mathbf{y} = f(\mathbf{t}) + \boldsymbol{\epsilon}$ with
$\boldsymbol{\epsilon} \sim \mathcal{N}(\mathbf{0}, \sigma_n^2 \mathbf{I})$,
the predictive distribution at test inputs $\mathbf{t}_*$ is:
\begin{align}
  \mathbf{f}_* | \mathbf{t}_*, \mathbf{t}, \mathbf{y}
    &\sim \mathcal{N}(\boldsymbol{\mu}_*, \boldsymbol{\Sigma}_*), \\
  \boldsymbol{\mu}_* &= m(\mathbf{t}_*) + \mathbf{K}_*^\top
    (\mathbf{K} + \sigma_n^2\mathbf{I})^{-1}(\mathbf{y} - m(\mathbf{t})),
    \label{eq:gp_mean} \\
  \boldsymbol{\Sigma}_* &= \mathbf{K}_{**} - \mathbf{K}_*^\top
    (\mathbf{K} + \sigma_n^2\mathbf{I})^{-1}\mathbf{K}_*,
    \label{eq:gp_var}
\end{align}
where $\mathbf{K} = k(\mathbf{t}, \mathbf{t})$,
$\mathbf{K}_* = k(\mathbf{t}, \mathbf{t}_*)$, and
$\mathbf{K}_{**} = k(\mathbf{t}_*, \mathbf{t}_*)$.

\subsubsection{Mat\'{e}rn Kernels}

The Mat\'{e}rn family of kernels provides control over the smoothness of
the GP sample paths via the parameter $\nu$:
\begin{equation}
  k_{\nu}(\Delta t) = \sigma_f^2 \frac{2^{1-\nu}}{\Gamma(\nu)}
    \left(\frac{\sqrt{2\nu}\,|\Delta t|}{\ell}\right)^\nu
    K_\nu\!\left(\frac{\sqrt{2\nu}\,|\Delta t|}{\ell}\right),
  \label{eq:matern}
\end{equation}
where $\ell$ is the length-scale, $\sigma_f^2$ the signal variance, and
$K_\nu$ the modified Bessel function. For $\nu = 5/2$, the kernel yields
twice-differentiable sample paths:
\begin{equation}
  k_{5/2}(\Delta t) = \sigma_f^2
    \left(1 + \frac{\sqrt{5}\,|\Delta t|}{\ell}
    + \frac{5\,\Delta t^2}{3\ell^2}\right)
    \exp\!\left(-\frac{\sqrt{5}\,|\Delta t|}{\ell}\right).
  \label{eq:matern52}
\end{equation}
This is our default temporal kernel for volume dynamics, encoding the prior
belief that tumour volume evolves smoothly over time.

\subsubsection{Multi-Output GPs and the Intrinsic Coregionalisation Model}

Multi-Output Gaussian Processes (MOGPs) extend scalar GPs to vector-valued
functions $\mathbf{f}: \R \to \R^D$, capturing correlations across output
dimensions. The Intrinsic Coregionalisation Model (ICM) \cite{alvarez2012kernels}
parameterises the cross-output covariance as a Kronecker product:
\begin{equation}
  \mathbf{K}(\mathbf{t}, \mathbf{t}') =
    \mathbf{B} \otimes k_t(\mathbf{t}, \mathbf{t}'),
  \label{eq:icm}
\end{equation}
where $\mathbf{B} \in \R^{D \times D}$ is a positive semi-definite
coregionalisation matrix encoding inter-output correlations, and
$k_t$ is a scalar temporal kernel shared across outputs. The Linear Model of
Coregionalisation (LMC) generalises this to a sum over $Q$ latent functions:
\begin{equation}
  \mathbf{K}(\mathbf{t}, \mathbf{t}') =
    \sum_{q=1}^{Q} \mathbf{B}_q \otimes k_q(\mathbf{t}, \mathbf{t}').
  \label{eq:lmc}
\end{equation}

% ---------------------------------------------------------------------------
% 2.5 Related Work
% ---------------------------------------------------------------------------
\subsection{Related Work}
\label{sec:bg:related}

\subsubsection{Brain Tumour Growth Prediction}

Tumour growth prediction from longitudinal MRI has been addressed through
several paradigms. Classical approaches fit parametric growth models
(exponential, Gompertz, logistic) to volume trajectories
\cite{nakasu2005serial,hashimoto2012natural}. Reaction--diffusion models
\cite{swanson2000quantitative,hormuth2021image} incorporate spatial diffusion
but require tissue-specific parameterisation and are computationally expensive.

Recent deep learning approaches include spatio-temporal convolutional LSTMs
\cite{zhang2020spatio} achieving 83.2\% Dice on glioma growth prediction,
GP-GAN \cite{elazab2020gp} combining Gaussian processes with GANs, and neural
fields conditioned on time \cite{chen2024vestibular} for vestibular schwannoma
growth. Puglisi et al.\ \cite{puglisi2025brain} demonstrated that a linear
mixed-effects model in a learned latent space produces clinically meaningful
progression predictions for brain atrophy, trained on 11,730~MRIs.

For meningiomas specifically, most existing work is limited to retrospective
volumetric analysis \cite{hashimoto2012natural,ius2020management,
nakamura2003natural}. To our knowledge, no prior work has applied foundation
model-derived latent spaces with GP-based temporal modelling to meningioma
growth prediction.

\subsubsection{Latent Space Temporal Prediction}

The Joint Embedding Predictive Architecture (JEPA) framework
\cite{lecun2022path} advocates predicting in representation space rather than
pixel space, avoiding the waste of capacity on irrelevant details. I-JEPA
\cite{assran2023self} and V-JEPA \cite{bardes2024revisiting} validated this
for images and video respectively. Longitudinal Self-Supervised Learning
(LSSL) \cite{zhao2021longitudinal} disentangles brain-age factors from latent
MRI representations on 811 ADNI subjects. Our pipeline aligns with this
``predict in latent space'' principle while using GPs instead of neural
predictors, providing closed-form uncertainty quantification for clinical
decision support.

\subsubsection{Domain Adaptation in Medical Imaging}

Domain adaptation for medical image analysis has been addressed through
adversarial training \cite{kamnitsas2017unsupervised}, style transfer
\cite{chen2022maxstyle}, and more recently through parameter-efficient
fine-tuning of foundation models \cite{dutt2023parameter}. Ben-David et al.\
\cite{bendavid2010theory} established the theoretical framework relating
domain divergence to target risk. Our work contributes empirical evidence on
LoRA adaptation for 3D medical imaging foundation models, including
dual-domain training that attempts to maintain glioma competence while
acquiring meningioma specialisation.


% --- Section 3: Data ---
% =============================================================================
% Section 3: Data
% =============================================================================
\section{Data}
\label{sec:data}

This section describes the three datasets, the unified HDF5 storage schema,
data splits, and the preprocessing pipeline.

% ---------------------------------------------------------------------------
% 3.1 Datasets
% ---------------------------------------------------------------------------
\subsection{Datasets}
\label{sec:data:datasets}

\subsubsection{BraTS-MEN 2024 (Cross-Sectional Meningioma)}

The BraTS Meningioma 2024 challenge dataset \cite{labella2024brats} comprises
1,000~subjects with pre-operative multi-parametric MRI (T1, T1ce, T2, FLAIR)
and expert manual segmentations. Each subject has a single study (timepoint).
Segmentation labels follow the BraTS convention: 0~(background), 1~(enhancing
tumour, ET), 2~(non-enhancing tumour, NET), and 3~(cystic component). This
dataset serves as the primary training source for Phases~1 and~2.

\subsubsection{BraTS-GLI 2021 (Longitudinal Glioma)}

The BraTS Glioma 2021 challenge dataset \cite{baid2021rsna} contains
1,350~scans from 613~patients (1--10 timepoints per patient) with the same
four MRI modalities. Glioma labels use: 0~(background), 1~(necrotic tumour
core, NETC), 2~(peritumoral oedema, SNFH), 3~(enhancing tumour, ET), and
4~(resection cavity, RC). This dataset is used for domain gap analysis
(\cref{sec:domain_gap}) and dual-domain LoRA training
(\cref{sec:lora:dual}).

\subsubsection{Andalusian Longitudinal Cohort (MenGrowth)}

A retrospective longitudinal cohort of 42~meningioma patients (137~total
studies) collected across hospitals in Andalusia, Spain. After quality control,
33~patients (100~studies) with 2--6~timepoints each are retained for growth
prediction. Observation statistics are summarised in \cref{tab:andalusian_stats}.

\begin{table}[H]
  \centering
  \caption{Andalusian longitudinal cohort statistics (after QC).}
  \label{tab:andalusian_stats}
  \begin{tabular}{lc}
    \toprule
    \textbf{Statistic} & \textbf{Value} \\
    \midrule
    Total patients ($N$) & 33 \\
    Total studies ($S$) & 100 \\
    Mean studies/patient ($\bar{n}$) & 3.03 \\
    Std studies/patient ($\sigma_n$) & 1.17 \\
    Range & [2, 6] \\
    Patients with exactly 2 studies & 19 (57.6\%) \\
    \bottomrule
  \end{tabular}
\end{table}

% ---------------------------------------------------------------------------
% 3.2 HDF5 Schema
% ---------------------------------------------------------------------------
\subsection{Unified HDF5 Schema}
\label{sec:data:h5}

All three datasets are stored in a unified HDF5 v2.0 schema, enabling a
single reader code path (\texttt{BraTSDatasetH5}) regardless of dataset type.
The schema structure is:

\begin{table}[H]
  \centering
  \caption{Unified HDF5 v2.0 schema. $N_s$: number of scans, $N_p$: number of
  patients.}
  \label{tab:h5_schema}
  \begin{tabular}{llp{8cm}}
    \toprule
    \textbf{Key} & \textbf{Shape/Type} & \textbf{Description} \\
    \midrule
    \texttt{attrs} & dict & $\{$\texttt{n\_scans}, \texttt{n\_patients},
      \texttt{roi\_size}, \texttt{spacing}, \texttt{channel\_order},
      \texttt{version}=``2.0'', \texttt{dataset\_type}, \texttt{domain}$\}$ \\
    \texttt{images} & $[N_s, 4, 192^3]$ float32 & Multi-modal MRI volumes \\
    \texttt{segs} & $[N_s, 1, 192^3]$ int8 & Segmentation masks \\
    \texttt{scan\_ids} & $[N_s]$ str & Unique scan identifiers \\
    \texttt{patient\_ids} & $[N_s]$ str & Patient identifiers \\
    \texttt{timepoint\_idx} & $[N_s]$ int32 & Per-patient timepoint index \\
    \texttt{semantic/} & group & Precomputed features: \texttt{volume} $[N_s, 4]$, \texttt{location} $[N_s, 3]$, \texttt{shape} $[N_s, 3]$ \\
    \texttt{longitudinal/} & group & CSR structure: \texttt{patient\_offsets} $[N_p+1]$, \texttt{patient\_list} $[N_p]$ \\
    \texttt{splits/} & group & Patient-level index arrays: \texttt{lora\_train}, \texttt{lora\_val}, \texttt{test} \\
    \texttt{metadata/} & group & \texttt{grade}, \texttt{age}, \texttt{sex} (when available) \\
    \bottomrule
  \end{tabular}
\end{table}

For cross-sectional data (BraTS-MEN), the longitudinal structure is trivial:
each patient has exactly one scan at timepoint~0, so $N_s = N_p$ and
\texttt{patient\_offsets} $= [0, 1, 2, \ldots, N]$. This design ensures that
all timepoints of a patient are assigned to the same split, preventing
patient-level data leakage.

% ---------------------------------------------------------------------------
% 3.3 Data Splits
% ---------------------------------------------------------------------------
\subsection{Data Splits}
\label{sec:data:splits}

Data splitting is performed at the patient level with seed 42 for
reproducibility. \Cref{tab:splits} summarises the splits for each dataset.

\begin{table}[H]
  \centering
  \caption{Data splits. BraTS-MEN splits are canonical; server configurations
    may differ slightly (see \cref{sec:lora:config}).}
  \label{tab:splits}
  \begin{tabular}{lccc}
    \toprule
    \textbf{Split} & \textbf{BraTS-MEN} & \textbf{BraTS-GLI (patients)} & \textbf{MenGrowth} \\
    \midrule
    \texttt{lora\_train} & 750 & 430 ($\sim$950 scans) & 23 \\
    \texttt{lora\_val}   & 100 & 60 ($\sim$130 scans)  & 3 \\
    \texttt{test}        & 150 & 123 ($\sim$270 scans) & 7 \\
    \midrule
    \textbf{Total} & 1,000 & 613 (1,350 scans) & 33 (100 scans) \\
    \bottomrule
  \end{tabular}
\end{table}

For Phase~2 (SDP), the canonical specification defines an additional
\texttt{sdp\_train} split of 225~subjects; however, the actual implementation
uses the \texttt{lora\_train} pool since the SDP operates on frozen encoder
features and there is no gradient contamination concern.

% ---------------------------------------------------------------------------
% 3.4 Preprocessing
% ---------------------------------------------------------------------------
\subsection{Preprocessing Pipeline}
\label{sec:data:preprocessing}

All volumes are spatially preprocessed to $192^3$ voxels at $1\,\text{mm}^3$
isotropic spacing during HDF5 conversion. Intensity normalisation is performed
at runtime. The critical channel ordering convention is
$[\text{FLAIR}, \text{T1ce}, \text{T1}, \text{T2}]$, matching
BrainSegFounder's training protocol---incorrect ordering causes near-zero
Dice ($\sim 0.00$) even on the glioma training domain.

\subsubsection{Training Transforms}

The training pipeline follows the BrainSegFounder protocol with minor
adaptations:
\begin{enumerate}[leftmargin=*]
  \item \texttt{CropForegroundd}: Crop to the non-zero brain region with
    \texttt{k\_divisible}$=[128, 128, 128]$.
  \item \texttt{SpatialPadd}: Pad to at least $128^3$.
  \item \texttt{RandSpatialCropd}: Random crop to $128^3$.
  \item \texttt{NormalizeIntensityd}: Per-channel, nonzero-voxel
    normalisation.
  \item Data augmentation: random flips (3 axes, $p=0.5$ each),
    intensity scaling (factor 0.1, $p=1.0$), intensity shift (offset 0.1,
    $p=1.0$).
\end{enumerate}

\subsubsection{Feature Extraction Transforms}

For feature extraction and evaluation, a larger $192^3$ ROI is used (via
\texttt{ResizeWithPadOrCropd}) to ensure 100\% tumour containment---the
$128^3$ centre crop captures only 38.8\% of meningioma tumours fully due
to their dural-based, peripheral location.

\subsubsection{Inference Transforms}

Full-resolution inference uses MONAI's sliding window inference with $128^3$
patches, 50\% overlap, and Gaussian blending.

% ---------------------------------------------------------------------------
% 3.5 Semantic Features
% ---------------------------------------------------------------------------
\subsection{Semantic Features}
\label{sec:data:semantic}

Semantic features are precomputed from segmentation masks and stored in the
HDF5 files. Three categories of features are extracted:

\begin{description}[leftmargin=*]
  \item[Volume (4 dims):] $[\log(V_{\text{total}}+1),\, \log(V_{\text{NCR}}+1),
    \, \log(V_{\text{ED}}+1),\, \log(V_{\text{ET}}+1)]$, where volumes are in
    mm$^3$. The log-transform stabilises the heavy-tailed distribution.

  \item[Location (3 dims):] Centroid coordinates $(c_z, c_y, c_x)$ in mm,
    computed as the centre of mass of the tumour segmentation mask.

  \item[Shape (3 dims):] Sphericity $\in (0, 1]$, $\log(\text{surface area})$,
    and solidity (ratio of tumour volume to convex hull volume), computed via
    \texttt{compute\_shape\_array()}.
\end{description}

\begin{remark}
  Following the R1 methodology revision (\cref{sec:eval:revision}), the
  volume target is reduced to a single whole-tumour log-volume
  $v^* = \log(V_{\text{WT}} + 1) \in \R^1$, and shape features are dropped
  from the SDP due to persistent low $R^2$ ($\leq 0.11$ across all
  conditions).
\end{remark}

% ---------------------------------------------------------------------------
% 3.6 Label Conversion
% ---------------------------------------------------------------------------
\subsection{Segmentation Label Conversion}
\label{sec:data:labels}

Both datasets are converted to a unified 3-channel hierarchical overlapping
representation for the segmentation loss:
\begin{align}
  \text{TC} &= \{x : \text{label}(x) \in S_{\text{TC}}\}, \\
  \text{WT} &= \{x : \text{label}(x) > 0\}, \\
  \text{ET} &= \{x : \text{label}(x) \in S_{\text{ET}}\},
\end{align}
where $S_{\text{TC}}$ and $S_{\text{ET}}$ are domain-dependent. For MEN:
$S_{\text{TC}} = \{1, 2\}$, $S_{\text{ET}} = \{1\}$. For GLI:
$S_{\text{TC}} = \{1, 3, 4\}$, $S_{\text{ET}} = \{3\}$. This conversion is
implemented in \texttt{segmentation.py.\_convert\_target()} and ensures that
the segmentation decoder output semantics are consistent across domains.


% --- Section 4: Domain Gap Analysis ---
% =============================================================================
% Section 4: Domain Gap Analysis
% =============================================================================
\section{Domain Gap Analysis}
\label{sec:domain_gap}

BrainSegFounder was pretrained and fine-tuned on glioma data. Before adapting
it to the meningioma domain, we quantify the distributional shift in the
frozen encoder's feature space to establish the empirical necessity of domain
adaptation.

% ---------------------------------------------------------------------------
% 4.1 Feature Extraction
% ---------------------------------------------------------------------------
\subsection{Feature Extraction Protocol}
\label{sec:dgap:extraction}

Features are extracted from the frozen BrainSegFounder encoder using the
\texttt{FeatureExtractor} pipeline:
\begin{equation}
  \mathbf{h} = \operatorname{GAP}\bigl(\texttt{encoder10}\bigl(
    \text{SwinViT}(\mathbf{x})\bigr)\bigr) \in \R^{768},
  \label{eq:feature_extraction}
\end{equation}
where $\mathbf{x} \in \R^{4 \times 192^3}$ is the input volume and
$\operatorname{GAP}$ denotes Global Average Pooling over the $4^3$ (or $6^3$
for $192^3$ input) spatial dimensions. We extract features from 200~random
BraTS-GLI subjects and 200~BraTS-MEN subjects from the test and SDP-train
splits.

% ---------------------------------------------------------------------------
% 4.2 Metrics
% ---------------------------------------------------------------------------
\subsection{Domain Gap Metrics}
\label{sec:dgap:metrics}

\subsubsection{Maximum Mean Discrepancy}

The Maximum Mean Discrepancy (MMD) \cite{gretton2012kernel} measures the
distance between two distributions in a reproducing kernel Hilbert space.
Given samples $\{\mathbf{h}_i^{\text{GLI}}\}_{i=1}^{n}$ and
$\{\mathbf{h}_j^{\text{MEN}}\}_{j=1}^{m}$, the unbiased estimator is:
\begin{equation}
  \widehat{\text{MMD}}^2 = \frac{1}{n(n-1)}\sum_{i \neq i'} k(\mathbf{h}_i,
    \mathbf{h}_{i'}) + \frac{1}{m(m-1)}\sum_{j \neq j'} k(\mathbf{h}_j,
    \mathbf{h}_{j'}) - \frac{2}{nm}\sum_{i,j} k(\mathbf{h}_i, \mathbf{h}_j),
  \label{eq:mmd}
\end{equation}
where $k$ is a Gaussian kernel with bandwidth $\sigma$ set to the median
pairwise distance (the median heuristic). Statistical significance is
assessed via a permutation test with $n_{\text{perm}} = 1{,}000$
permutations.

\subsubsection{Centred Kernel Alignment}

CKA \cite{kornblith2019similarity} measures representational similarity
between two feature matrices $\mathbf{H}_1 \in \R^{n \times d_1}$ and
$\mathbf{H}_2 \in \R^{n \times d_2}$ via:
\begin{equation}
  \text{CKA}(\mathbf{H}_1, \mathbf{H}_2) = \frac{
    \|\bar{\mathbf{H}}_1^\top \bar{\mathbf{H}}_2\|_F^2
  }{
    \|\bar{\mathbf{H}}_1^\top \bar{\mathbf{H}}_1\|_F \cdot
    \|\bar{\mathbf{H}}_2^\top \bar{\mathbf{H}}_2\|_F
  },
  \label{eq:cka}
\end{equation}
where $\bar{\mathbf{H}}$ denotes column-centred matrices. CKA $\in [0, 1]$,
with 1 indicating identical representations (up to linear transformation).

\subsubsection{Linear Probes}

Following Alain and Bengio \cite{alain2017understanding}, we train Ridge
regression probes from frozen features $\mathbf{h}$ to semantic targets
(volume, location, shape). Cross-validated $R^2$ quantifies the linear
decodability of each semantic factor. Low $R^2$ on the target domain indicates
that the frozen encoder has not learned meningioma-relevant features.

% ---------------------------------------------------------------------------
% 4.3 Results
% ---------------------------------------------------------------------------
\subsection{Results}
\label{sec:dgap:results}

\begin{table}[H]
  \centering
  \caption{Domain gap metrics between frozen BrainSegFounder features on
    BraTS-GLI vs.\ BraTS-MEN. Results confirm significant distributional shift
    motivating domain adaptation.}
  \label{tab:domain_gap}
  \begin{tabular}{lcc}
    \toprule
    \textbf{Metric} & \textbf{Value} & \textbf{Interpretation} \\
    \midrule
    $\text{MMD}^2$ (GLI vs.\ MEN) & \textit{pending} & $> 0$ confirms shift \\
    MMD $p$-value (permutation test) & \textit{pending} & $< 0.05$ confirms significance \\
    CKA (GLI vs.\ MEN) & \textit{pending} & Low $\Rightarrow$ different structure \\
    \midrule
    Linear probe $R^2$ (volume, MEN) & \textit{pending} & Baseline decodability \\
    Linear probe $R^2$ (location, MEN) & \textit{pending} & Baseline decodability \\
    Linear probe $R^2$ (shape, MEN) & \textit{pending} & Baseline decodability \\
    \bottomrule
  \end{tabular}
\end{table}

\begin{figure}[H]
  \centering
  % \includegraphics[width=0.8\textwidth]{figures/domain_umap.pdf}
  \fbox{\parbox{0.8\textwidth}{\centering\vspace{3cm}
    \textbf{Figure placeholder:} UMAP of frozen BrainSegFounder features.\\
    GLI (orange) vs.\ MEN (blue). Expected: clear separation.
  \vspace{3cm}}}
  \caption{UMAP visualisation of frozen BrainSegFounder encoder features
    ($\mathbf{h} \in \R^{768}$) for BraTS-GLI (orange, $n=200$) and BraTS-MEN
    (blue, $n=200$). The visible separation confirms the domain gap motivating
    LoRA adaptation.}
  \label{fig:domain_umap}
\end{figure}

The domain gap analysis confirms that the frozen glioma-trained encoder
produces features with significant distributional shift when applied to
meningioma data. The MMD permutation test rejects the null hypothesis of
identical distributions, and the linear probe $R^2$ values on MEN data are
substantially lower than within-domain GLI probes, indicating that the frozen
features inadequately capture meningioma-specific semantic content. These
findings motivate the LoRA adaptation described in \cref{sec:lora}.

% ---------------------------------------------------------------------------
% 4.4 Effective Rank Analysis
% ---------------------------------------------------------------------------
\subsection{Effective Rank Analysis}
\label{sec:dgap:rank}

The effective rank of the feature matrix, defined as the exponential of the
spectral entropy:
\begin{equation}
  \text{erank}(\mathbf{H}) = \exp\!\left(-\sum_{i=1}^{d} \bar{s}_i
    \log \bar{s}_i\right),
  \quad \bar{s}_i = \frac{s_i}{\sum_j s_j},
  \label{eq:erank}
\end{equation}
where $s_i$ are the singular values of $\mathbf{H}$, measures the
distributional uniformity of the representation. The frozen BrainSegFounder
encoder achieves an effective rank of approximately 47.6 on MEN data,
indicating moderate utilisation of the 768-dimensional feature space.


% --- Section 5: Phase 1 — LoRA Encoder Adaptation ---
% =============================================================================
% Section 5: Phase 1 — LoRA Encoder Adaptation
% =============================================================================
\section{Phase 1: LoRA Encoder Adaptation}
\label{sec:lora}

% ---------------------------------------------------------------------------
% 5.1 Adaptation Strategy
% ---------------------------------------------------------------------------
\subsection{Adaptation Strategy}
\label{sec:lora:strategy}

We adapt the frozen BrainSegFounder encoder to the meningioma domain via LoRA
(\cref{sec:bg:lora}), injecting low-rank adapter matrices into the
self-attention projections of Stages~3--4. Stages~0--2 remain frozen to
preserve low-level anatomical features learned from the UK~Biobank SSL
pretraining. Segmentation on BraTS-MEN 2024 serves as the proxy task.

\subsubsection{Target Modules}

LoRA adapters are injected into the $\mathbf{Q}$, $\mathbf{K}$, $\mathbf{V}$
projection matrices within the multi-head self-attention blocks of Stages~3
and~4:
\begin{align}
  &\texttt{swinViT.layers3.0.blocks.\{0,1\}.attn.qkv.weight}, \\
  &\texttt{swinViT.layers4.0.blocks.\{0,1\}.attn.qkv.weight}.
\end{align}
This yields 4 adapted layers, each with rank-$r$ decompositions for Q, K,
and V (combined as a single QKV projection in the Swin implementation). The
total number of trainable LoRA parameters is approximately $197\text{K}$ for
$r=8$, compared to 62M total model parameters ($<0.32\%$).

\subsubsection{LoRA Hyperparameters}

\begin{table}[H]
  \centering
  \caption{LoRA configuration (selected via rank ablation, see
    \cref{sec:lora:ablation}).}
  \label{tab:lora_config}
  \begin{tabular}{lcp{7cm}}
    \toprule
    \textbf{Parameter} & \textbf{Value} & \textbf{Rationale} \\
    \midrule
    Rank ($r$) & 8 & Optimal from ablation over $r \in \{2, 4, 8, 16, 32\}$ \\
    Scale ($\alpha$) & 16 & Effective scale $\alpha/r = 2.0$ \\
    Dropout & 0.1 & Regularisation \\
    Target stages & 3, 4 & High-level features; stages 0--2 frozen \\
    Target modules & QKV & All attention projections \\
    \bottomrule
  \end{tabular}
\end{table}

% ---------------------------------------------------------------------------
% 5.2 Auxiliary Semantic Heads
% ---------------------------------------------------------------------------
\subsection{Auxiliary Semantic Heads}
\label{sec:lora:aux}

To enrich the encoder's feature space for downstream SDP training, we
attach auxiliary semantic regression heads that predict tumour properties
directly from GAP-pooled encoder features:
\begin{equation}
  \hat{\mathbf{y}}_p = \pi_p\bigl(\operatorname{GAP}(\texttt{encoder10}(\mathbf{x}))\bigr),
  \quad p \in \{\text{vol}\},
  \label{eq:aux_heads}
\end{equation}
where $\pi_{\text{vol}}: \R^{768} \to \R^{1}$ is a linear head predicting
log whole-tumour volume. The auxiliary loss is:
\begin{equation}
  \mathcal{L}_{\text{aux}} = \lambda_{\text{aux}} \cdot
    \|\hat{\mathbf{y}}_{\text{vol}} - \mathbf{y}_{\text{vol}}\|_2^2.
  \label{eq:aux_loss}
\end{equation}
Following the R1 revision (\cref{sec:eval:revision}), auxiliary heads for
location and shape have been removed; only the volume head is retained. The
auxiliary loss is activated after epoch~5 with a linear warmup over 10~epochs
($\lambda_{\text{aux}} = 0.1$ at full strength).

\begin{remark}
  The auxiliary heads are discarded after Phase~1. Their sole purpose is to
  provide multi-task learning signal during encoder adaptation, encouraging
  the encoder to produce features from which tumour volume can be linearly
  decoded.
\end{remark}

% ---------------------------------------------------------------------------
% 5.3 Training Protocol
% ---------------------------------------------------------------------------
\subsection{Training Protocol}
\label{sec:lora:training}

\begin{table}[H]
  \centering
  \caption{Phase 1 training configuration.}
  \label{tab:lora_training}
  \begin{tabular}{lcp{6cm}}
    \toprule
    \textbf{Parameter} & \textbf{Value} & \textbf{Rationale} \\
    \midrule
    Epochs & 150 & With early stopping (patience 25) \\
    LR (LoRA) & $10^{-4}$ & Standard for LoRA adaptation \\
    LR (decoder) & $5 \times 10^{-4}$ & Larger for pretrained decoder \\
    LR (aux heads) & $10^{-3}$ & Small heads, fast convergence \\
    Optimizer & AdamW & Decoupled weight decay \\
    Weight decay & $10^{-5}$ & Light regularisation \\
    Scheduler & Cosine decay, 10-epoch warmup & Smooth convergence \\
    Batch size & 2 (per GPU) & Memory-constrained by $128^3$ input \\
    Gradient accumulation & 4 steps & Effective batch size 8 \\
    Gradient clipping & $\|\mathbf{g}\|_2 \leq 1.0$ & Stability \\
    Precision & bf16-mixed & Memory efficiency on A100 \\
    \bottomrule
  \end{tabular}
\end{table}

The segmentation loss combines Dice loss and cross-entropy with equal weight:
\begin{equation}
  \mathcal{L}_{\text{seg}} = \lambda_{\text{dice}}\,\mathcal{L}_{\text{Dice}}
    + \lambda_{\text{CE}}\,\mathcal{L}_{\text{CE}},
  \quad \lambda_{\text{dice}} = \lambda_{\text{CE}} = 1.0.
  \label{eq:seg_loss}
\end{equation}
The total training objective is:
\begin{equation}
  \mathcal{L} = \mathcal{L}_{\text{seg}} + w(e)\,\mathcal{L}_{\text{aux}},
  \label{eq:lora_total_loss}
\end{equation}
where $w(e) = \lambda_{\text{aux}} \cdot \min(1, \max(0, (e - 5)/10))$ is
the warmup schedule for the auxiliary loss starting at epoch~5.

% ---------------------------------------------------------------------------
% 5.4 Post-Training Protocol
% ---------------------------------------------------------------------------
\subsection{Post-Training: LoRA Merge and Checkpoint Export}
\label{sec:lora:merge}

After training, LoRA weights are absorbed into the base model weights:
\begin{equation}
  \mathbf{W}_{\text{merged}} = \mathbf{W}_0 + \frac{\alpha}{r}\,
    \mathbf{B}\mathbf{A}.
  \label{eq:lora_merge}
\end{equation}
The segmentation decoder and auxiliary heads are then discarded, yielding a
merged encoder checkpoint (\texttt{phase1\_encoder\_merged.pt}) containing
only the SwinViT and encoder1--10 weights. All encoder parameters are frozen
for subsequent phases.

% ---------------------------------------------------------------------------
% 5.5 Ablation Study
% ---------------------------------------------------------------------------
\subsection{Ablation Results}
\label{sec:lora:ablation}

\subsubsection{Rank Sweep (A1)}

We evaluate LoRA ranks $r \in \{2, 4, 8, 16, 32\}$ on segmentation Dice
and linear probe $R^2$ for volume prediction.

\begin{table}[H]
  \centering
  \caption{LoRA rank ablation (A1). Dice scores on \texttt{lora\_val} and
    linear probe $R^2$ for volume on test features.}
  \label{tab:rank_ablation}
  \begin{tabular}{cccccc}
    \toprule
    \textbf{Rank} & \textbf{LoRA params} & \textbf{Dice (WT)} &
      \textbf{Dice (TC)} & \textbf{Dice (ET)} & \textbf{Vol $R^2$} \\
    \midrule
    Frozen & 0 & \textit{pending} & \textit{pending} & \textit{pending} & \textit{pending} \\
    2  & $\sim$49K  & \textit{pending} & \textit{pending} & \textit{pending} & \textit{pending} \\
    4  & $\sim$99K  & \textit{pending} & \textit{pending} & \textit{pending} & \textit{pending} \\
    8  & $\sim$197K & \textit{pending} & \textit{pending} & \textit{pending} & \textit{pending} \\
    16 & $\sim$394K & \textit{pending} & \textit{pending} & \textit{pending} & \textit{pending} \\
    32 & $\sim$787K & \textit{pending} & \textit{pending} & \textit{pending} & \textit{pending} \\
    \bottomrule
  \end{tabular}
\end{table}

\begin{figure}[H]
  \centering
  \fbox{\parbox{0.8\textwidth}{\centering\vspace{3cm}
    \textbf{Figure placeholder:} LoRA rank ablation curves.\\
    Left: Dice vs.\ rank. Right: Probe $R^2$ vs.\ rank.
  \vspace{3cm}}}
  \caption{LoRA rank ablation (A1). Left panel shows validation Dice scores
    per sub-region as a function of LoRA rank. Right panel shows linear probe
    $R^2$ for volume prediction. Rank $r=8$ achieves the best trade-off
    between capacity and parameter efficiency.}
  \label{fig:rank_ablation}
\end{figure}

\subsubsection{LoRA vs.\ DoRA (A2)}

DoRA (\cref{eq:dora}) is evaluated as an alternative adapter type at $r=8$.

\begin{table}[H]
  \centering
  \caption{LoRA vs.\ DoRA comparison (A2) at rank 8.}
  \label{tab:lora_vs_dora}
  \begin{tabular}{lccc}
    \toprule
    \textbf{Adapter} & \textbf{Dice (WT)} & \textbf{Vol $R^2$ (linear)} & \textbf{Vol $R^2$ (RBF)} \\
    \midrule
    LoRA ($r=8$) & \textit{pending} & \textit{pending} & \textit{pending} \\
    DoRA ($r=8$) & \textit{pending} & \textit{pending} & \textit{pending} \\
    \bottomrule
  \end{tabular}
\end{table}

% ---------------------------------------------------------------------------
% 5.6 Dual-Domain Adaptation
% ---------------------------------------------------------------------------
\subsection{Dual-Domain Adaptation}
\label{sec:lora:dual}

A key scientific question is whether a single LoRA-adapted encoder can produce
semantically structured features for both glioma and meningioma simultaneously.
Dual-domain training combines BraTS-MEN ($\sim$750~studies) and BraTS-GLI
($\sim$950~scans) via a \texttt{ConcatDataset} with domain-balanced
\texttt{WeightedRandomSampler} (50/50 domain weight).

\subsubsection{Empirical Findings}

\begin{table}[H]
  \centering
  \caption{Single-domain vs.\ dual-domain LoRA adaptation results.}
  \label{tab:dual_domain}
  \begin{tabular}{lcccccc}
    \toprule
    & \textbf{MEN Dice} & \textbf{Vol $R^2$} & \textbf{CV Vol $R^2$}
    & \textbf{MMD$^2$} & \textbf{CKA} & \textbf{erank} \\
    \midrule
    Frozen baseline & 0.52 & 0.56 & --- & 0.118 & 0.002 & 47.6 \\
    MEN LoRA $r$=8  & 0.48 & 0.57 & --- & 0.116 & 0.002 & 45.8 \\
    Dual LoRA $r$=8 & 0.57 & 0.62 & 0.712 & 0.037 & 0.017 & 25.4 \\
    \bottomrule
  \end{tabular}
\end{table}

The dual-domain encoder achieves the best volume probe $R^2$ and reduces the
domain gap (MMD$^2$: $0.118 \to 0.037$, a 68.6\% reduction), but at the cost
of a substantial drop in effective rank ($47.6 \to 25.4$, a 46.6\% reduction),
indicating representation compression.

\subsubsection{Critical Finding: Cross-Domain Transfer Failure}

Despite the reduced domain gap metrics, cross-domain GP probes reveal a
critical limitation: training probes on GLI features and evaluating on MEN
features yields $R^2 \approx -0.01$ for volume, indicating that the shared
feature space does not support cross-domain semantic transfer. Furthermore,
MEN segmentation Dice degrades from 0.64 to 0.36 under dual training,
indicating catastrophic interference rather than beneficial transfer.

\begin{remark}
  The dual-domain results provide important negative evidence: domain gap
  reduction (measured by MMD, CKA) does not guarantee cross-domain semantic
  compatibility. The effective rank drop from 47.6 to 25.2 suggests that the
  encoder resolves the domain conflict by compressing both domains into a
  lower-dimensional shared subspace that loses domain-specific semantic
  content.
\end{remark}

\subsubsection{LoRA Marginality}

A surprising finding is the marginal benefit of LoRA adaptation for semantic
probes: the frozen baseline achieves volume $R^2 = 0.794$, while the best
LoRA condition improves this by only $\Delta R^2 = +0.016$ (within noise).
This suggests that BrainSegFounder's SSL pretraining already produces
features with reasonable meningioma volume decodability, and the primary
benefit of LoRA is in segmentation quality rather than downstream feature
quality.

% ---------------------------------------------------------------------------
% 5.7 Configuration
% ---------------------------------------------------------------------------
\subsection{Configuration}
\label{sec:lora:config}

All LoRA experiments are configured via YAML files under
\texttt{experiments/lora/config/}. Server configurations use 250/50/200/500
splits (lora\_train/lora\_val/probe\_train/test), differing slightly from the
canonical specification. Four conditions are defined:
LoRA/DoRA $\times$ semantic/no-semantic heads.


% --- Section 6: Phase 2 — Supervised Disentangled Projection ---
% =============================================================================
% Section 6: Phase 2 — Supervised Disentangled Projection
% =============================================================================
\section{Phase 2: Supervised Disentangled Projection}
\label{sec:sdp}

% ---------------------------------------------------------------------------
% 6.1 Motivation
% ---------------------------------------------------------------------------
\subsection{Motivation}
\label{sec:sdp:motivation}

The encoder features $\mathbf{h} \in \R^{768}$ are high-dimensional and
unstructured: semantic factors (volume, location, shape) are entangled across
all 768~dimensions. Downstream GP-based growth prediction requires features
that are (i)~\emph{semantically interpretable}: tumour volume should be
decodable from a known subset of latent dimensions; (ii)~\emph{disentangled}:
volume, location, and shape should occupy independent subspaces;
(iii)~\emph{low-dimensional}: the GP covariance matrix scales as $O(D^2)$ in
the output dimensionality $D$.

The Supervised Disentangled Projection (SDP) addresses all three requirements
through a lightweight MLP trained on frozen encoder features with a composite
loss combining semantic supervision, collapse prevention, and independence
regularisation.

\subsubsection{Why Not a VAE?}

Early experiments with a Semi-Supervised VAE (SemiVAE) were abandoned due to
three failure modes: (i)~posterior collapse, where latent dimensions become
uninformative despite extensive hyperparameter tuning (7-phase curriculum
across 600~epochs); (ii)~KL distortion, where the $D_{\KL}(q \| p)$ penalty
actively minimises the mutual information $I(x; z)$ between input and
representation, destroying the rich foundation model features; and
(iii)~reconstruction burden, where the decoder wastes capacity on
pixel-level fidelity irrelevant to downstream growth prediction.

The SDP replaces the generative objective with a discriminative one that
directly optimises for the properties needed by the GP temporal model,
following the theoretical mandate of Locatello et al.\
\cite{locatello2019challenging}: supervised disentanglement resolves the
impossibility of unsupervised approaches.

% ---------------------------------------------------------------------------
% 6.2 Architecture
% ---------------------------------------------------------------------------
\subsection{Architecture}
\label{sec:sdp:architecture}

The SDP consists of a projection network $g_\theta$ and per-partition semantic
heads $\{\pi_p\}$:
\begin{equation}
  g_\theta: \R^{768} \to \R^{128}, \quad
  g_\theta(\mathbf{h}) = \text{SN}(\mathbf{W}_2)\,
    \sigma\bigl(\text{SN}(\mathbf{W}_1)\,
    \text{LN}(\mathbf{h}) + \mathbf{b}_1\bigr) + \mathbf{b}_2,
  \label{eq:sdp_forward}
\end{equation}
where $\text{LN}$ is LayerNorm, $\sigma$ is GELU with dropout $p=0.1$, and
$\text{SN}(\cdot)$ denotes spectral normalisation applied to \emph{both}
linear layers (\cref{sec:bg:disentanglement}). The architecture is:

\begin{equation}
  \underbrace{\mathbf{h} \in \R^{768}}_{\text{frozen encoder}}
  \xrightarrow{\text{LN}}
  \xrightarrow{\text{SN-Linear}(768, 512)}
  \xrightarrow{\text{GELU + Drop}}
  \xrightarrow{\text{SN-Linear}(512, 128)}
  \underbrace{\mathbf{z} \in \R^{128}}_{\text{structured latent}}
  \label{eq:sdp_pipeline}
\end{equation}

The total parameter count is approximately 459K
($768 \times 512 + 512 + 512 \times 128 + 128$).

% ---------------------------------------------------------------------------
% 6.3 Partition Structure
% ---------------------------------------------------------------------------
\subsection{Latent Partition Structure}
\label{sec:sdp:partitions}

The 128-dimensional latent vector is partitioned into semantically supervised
and unsupervised subspaces:

\begin{table}[H]
  \centering
  \caption{Original SDP partition structure ($3+1$ partitions). See
    \cref{sec:eval:revision} for the R1-revised $1+1$ structure.}
  \label{tab:partitions}
  \begin{tabular}{lcccc}
    \toprule
    \textbf{Partition} & \textbf{Dims} & \textbf{Indices} &
      \textbf{Target dim} & \textbf{Semantic head} \\
    \midrule
    $\Vol$ (volume)   & 24 & $[0, 24)$    & 4 & $\pi_{\text{vol}}: \R^{24} \to \R^4$ \\
    $\Loc$ (location) & 8  & $[24, 32)$   & 3 & $\pi_{\text{loc}}: \R^{8} \to \R^3$ \\
    $\Shape$ (shape)  & 12 & $[32, 44)$   & 6 & $\pi_{\text{shape}}: \R^{12} \to \R^6$ \\
    $\Res$ (residual) & 84 & $[44, 128)$  & --- & None (unsupervised) \\
    \midrule
    \textbf{Total}    & 128 & $[0, 128)$ & 13 & --- \\
    \bottomrule
  \end{tabular}
\end{table}

Each supervised partition $p$ has a linear semantic head
$\pi_p: \R^{d_p} \to \R^{k_p}$ that predicts ground-truth semantic features
extracted from segmentation masks. The residual partition captures
information not explained by the three supervised factors.

\begin{remark}[R1 Revision]
  Following empirical results (shape $R^2 \leq 0.11$, temporally static
  meningioma location), the partition structure is revised to
  $\mathbf{z} = [\Vol \in \R^{32} \mid \Res \in \R^{96}]$ with a single
  volume head $\pi_{\text{vol}}: \R^{32} \to \R^1$ predicting
  $\log(V_{\text{WT}} + 1)$. Location is relocated to a static GP covariate
  (\cref{sec:eval:revision}).
\end{remark}

% ---------------------------------------------------------------------------
% 6.4 Loss Function
% ---------------------------------------------------------------------------
\subsection{Composite Loss Function}
\label{sec:sdp:loss}

The SDP training objective combines four terms:
\begin{equation}
  \mathcal{L}_{\SDP} =
    \underbrace{\sum_{p \in \mathcal{P}} \lambda_p\,
      \|\hat{\mathbf{y}}_p - \mathbf{y}_p\|_2^2
    }_{\text{Semantic regression}}
  + \underbrace{\lambda_{\text{var}}\,
      \mathcal{L}_{\text{var}}(\mathbf{z})
    }_{\text{Collapse prevention}}
  + \underbrace{\lambda_{\text{cov}}\,
      \mathcal{L}_{\text{cov}}(\mathbf{z})
    }_{\text{Linear independence}}
  + \underbrace{\lambda_{\text{dCor}}\!
      \sum_{\substack{i,j \in \mathcal{P} \\ i < j}}
      \dCor(\mathbf{z}_i, \mathbf{z}_j)
    }_{\text{Nonlinear independence}},
  \label{eq:sdp_loss}
\end{equation}
where $\mathcal{P}$ is the set of supervised partitions,
$\hat{\mathbf{y}}_p = \pi_p(\mathbf{z}_p)$, and each loss term serves a
distinct purpose.

\subsubsection{Semantic Regression}

For each supervised partition, MSE loss anchors the partition's information
content to a known semantic factor:
\begin{equation}
  \mathcal{L}_p^{\text{sem}} = \frac{1}{k_p}\,
    \|\pi_p(\mathbf{z}_p) - \mathbf{y}_p\|_2^2.
  \label{eq:sem_loss}
\end{equation}
Targets $\mathbf{y}_p$ are normalised to zero mean and unit variance using
statistics computed on the training pool only (\cref{sec:sdp:normalisation}).

\subsubsection{Variance Hinge Loss (Collapse Prevention)}

The VICReg variance term (\cref{eq:vicreg_var}) is applied to all 128
dimensions of $\mathbf{z}$, ensuring no dimension collapses to a constant.
The threshold $\gamma = 1.0$ targets unit standard deviation per dimension.

\subsubsection{Covariance Penalty (Linear Independence)}

Cross-partition covariance is penalised:
\begin{equation}
  \mathcal{L}_{\text{cov}} = \frac{1}{d^2}
    \sum_{\substack{i,j: \\ \text{part}(i) \neq \text{part}(j)}} C_{ij}^2,
  \label{eq:cov_loss}
\end{equation}
where $\mathbf{C}$ is the batch covariance matrix and $\text{part}(i)$
denotes the partition containing dimension $i$. Only cross-partition
off-diagonal elements are penalised; within-partition correlations are
permitted.

\subsubsection{Distance Correlation (Nonlinear Independence)}

For each pair of supervised partitions $(i, j)$, the distance correlation
$\dCor(\mathbf{z}_i, \mathbf{z}_j)$ (\cref{eq:dcor}) captures nonlinear
dependencies not detected by the covariance penalty. The empirical distance
correlation is computed from $O(N^2)$ pairwise distance matrices, which for
$N = 800$ (full-batch training) requires approximately 30~MB of GPU memory
per partition pair.

% ---------------------------------------------------------------------------
% 6.5 Curriculum Training
% ---------------------------------------------------------------------------
\subsection{Curriculum Training Schedule}
\label{sec:sdp:curriculum}

Following the capacity-annealing principle from the disentanglement literature
\cite{burgess2018understanding}, losses are activated in a staged curriculum:

\begin{table}[H]
  \centering
  \caption{SDP curriculum training schedule (4 phases).}
  \label{tab:curriculum}
  \begin{tabular}{clp{7cm}}
    \toprule
    \textbf{Phase} & \textbf{Epochs} & \textbf{Active losses} \\
    \midrule
    Warm-up     & 0--9    & $\mathcal{L}_{\text{var}}$ only \\
    Semantic    & 10--39  & $+ \mathcal{L}_{\text{vol}}, \mathcal{L}_{\text{loc}}, \mathcal{L}_{\text{shape}}$ \\
    Independence & 40--59 & $+ \mathcal{L}_{\text{cov}}, \mathcal{L}_{\text{dCor}}$ \\
    Full         & 60--100 & All losses at full strength \\
    \bottomrule
  \end{tabular}
\end{table}

The rationale is that semantic regression must first establish the partition
structure before independence regularisation is imposed. Premature
independence penalties can collapse dimensions before they have acquired
semantic content.

% ---------------------------------------------------------------------------
% 6.6 Hyperparameters
% ---------------------------------------------------------------------------
\subsection{Hyperparameters}
\label{sec:sdp:hyperparams}

\begin{table}[H]
  \centering
  \caption{SDP training hyperparameters.}
  \label{tab:sdp_hyperparams}
  \begin{tabular}{lcp{6cm}}
    \toprule
    \textbf{Parameter} & \textbf{Value} & \textbf{Rationale} \\
    \midrule
    $\lambda_{\text{vol}}$ & 20.0 & Highest weight; volume is primary target \\
    $\lambda_{\text{loc}}$ & 12.0 & Moderate weight \\
    $\lambda_{\text{shape}}$ & 15.0 & Higher than location (harder to learn) \\
    $\lambda_{\text{cov}}$ & 5.0 & Cross-partition linear independence \\
    $\lambda_{\text{var}}$ & 5.0 & Collapse prevention \\
    $\lambda_{\text{dCor}}$ & 2.0 & Nonlinear independence \\
    Learning rate & $10^{-3}$ & Standard for small MLPs \\
    Optimizer & AdamW & Decoupled weight decay \\
    Weight decay & 0.01 & Regularisation \\
    Epochs & 100 & Sufficient for convergence \\
    Batch size & 800 (full-batch) & Exact covariance/dCor estimates \\
    Scheduler & Cosine, 5-epoch warmup & Smooth convergence \\
    \bottomrule
  \end{tabular}
\end{table}

The full-batch training is computationally trivial since the SDP operates on
precomputed 768-dimensional feature vectors, not 3D volumes. Full-batch
provides exact (not estimated) covariance and distance correlation statistics.

% ---------------------------------------------------------------------------
% 6.7 Normalisation Scope
% ---------------------------------------------------------------------------
\subsection{Normalisation Scope}
\label{sec:sdp:normalisation}

Target normalisation statistics ($\mu$, $\sigma$) are computed on the training
pool (800~subjects) only and applied without recomputation to validation
(100), test (150), and the Andalusian cohort. This prevents information
leakage and ensures consistent scaling across all downstream uses of the
semantic heads.

% ---------------------------------------------------------------------------
% 6.8 Lipschitz Analysis
% ---------------------------------------------------------------------------
\subsection{Lipschitz Continuity Analysis}
\label{sec:sdp:lipschitz}

With spectral normalisation on both linear layers and GELU activation
(Lipschitz constant $\approx 1.13$), the global Lipschitz constant of the SDP
is bounded by:
\begin{equation}
  \operatorname{Lip}(g_\theta) \leq
    \underbrace{1}_{\text{LN}} \cdot
    \underbrace{1}_{\text{SN}(\mathbf{W}_1)} \cdot
    \underbrace{1.13}_{\text{GELU}} \cdot
    \underbrace{1}_{\text{SN}(\mathbf{W}_2)} \approx 1.13.
  \label{eq:lipschitz}
\end{equation}
This ensures that the projection is a near-contraction:
$\|\mathbf{z}_1 - \mathbf{z}_2\| \leq 1.13\,\|\mathbf{h}_1 - \mathbf{h}_2\|$.
Encoder feature noise is not amplified through the projection, yielding
well-conditioned inputs for downstream GP models.

% ---------------------------------------------------------------------------
% 6.9 Quality Targets
% ---------------------------------------------------------------------------
\subsection{Quality Targets}
\label{sec:sdp:targets}

\begin{table}[H]
  \centering
  \caption{SDP quality targets (validation set).}
  \label{tab:sdp_targets}
  \begin{tabular}{lccc}
    \toprule
    \textbf{Metric} & \textbf{Minimum} & \textbf{Target} & \textbf{Blocking?} \\
    \midrule
    Volume $R^2$ & 0.80 & 0.90 & Yes \\
    Location $R^2$ & 0.85 & 0.95 & Yes \\
    Shape $R^2$ & 0.30 & 0.40 & Yes \\
    Max cross-partition correlation & $< 0.30$ & $< 0.20$ & Yes \\
    $\dCor(\text{vol}, \text{loc})$ & $< 0.20$ & $< 0.10$ & Yes \\
    Dims with variance $> 0.5$ & $\geq 85\%$ & $\geq 95\%$ & Yes \\
    \bottomrule
  \end{tabular}
\end{table}

\begin{remark}
  The shape $R^2$ target of 0.30 was not achieved in any experimental
  condition ($R^2 \leq 0.11$), motivating the R1 revision to drop the shape
  partition entirely (\cref{sec:eval:revision}).
\end{remark}


% --- Section 7: Phase 3 — Cohort Encoding and Harmonisation ---
% =============================================================================
% Section 7: Phase 3 — Cohort Encoding and Harmonisation
% =============================================================================
\section{Phase 3: Cohort Encoding and Harmonisation}
\label{sec:encoding}

% ---------------------------------------------------------------------------
% 7.1 Overview
% ---------------------------------------------------------------------------
\subsection{Overview}
\label{sec:enc:overview}

Phase~3 encodes all MRI volumes from both the BraTS-MEN cross-sectional
cohort and the Andalusian longitudinal cohort through the frozen
encoder--SDP pipeline, producing per-subject latent vectors and per-patient
temporal trajectories suitable for GP-based growth prediction.

The pipeline is fully deterministic (no dropout, model in evaluation mode)
and proceeds as:
\begin{equation}
  \mathbf{x} \in \R^{4 \times 192^3}
  \xrightarrow{\text{SwinViT (frozen)}}
  \xrightarrow{\texttt{encoder10}}
  \xrightarrow{\text{GAP}}
  \mathbf{h} \in \R^{768}
  \xrightarrow{g_\theta\,\text{(frozen SDP)}}
  \mathbf{z} \in \R^{128}.
  \label{eq:encoding_pipeline}
\end{equation}

% ---------------------------------------------------------------------------
% 7.2 Sliding Window Feature Extraction
% ---------------------------------------------------------------------------
\subsection{Sliding Window Feature Extraction}
\label{sec:enc:sliding_window}

For volumes larger than the $128^3$ training ROI, or to ensure consistent
feature extraction regardless of tumour position, we employ a sliding window
approach with tumour-weighted averaging:
\begin{equation}
  \mathbf{h} = \frac{\sum_{p=1}^{P} w_p \cdot \operatorname{GAP}
    \bigl(\texttt{encoder10}(\mathbf{x}_p)\bigr)}
    {\sum_{p=1}^{P} w_p},
  \quad w_p = \frac{|\text{mask}_p \cap \text{tumour}|}{|\text{patch}_p|},
  \label{eq:sliding_window}
\end{equation}
where $\mathbf{x}_p$ denotes patch $p$, and $w_p$ weights each patch's
contribution by its tumour content fraction. Patches with no tumour
($w_p = 0$) are excluded. This ensures that the feature vector is dominated
by tumour-containing regions.

In practice, the feature extraction transforms use the full $192^3$ ROI
(see \cref{sec:data:preprocessing}), which provides 100\% tumour containment
for meningiomas without requiring sliding window aggregation. The sliding
window protocol is retained for robustness and for use with the Andalusian
cohort where acquisition protocols may vary.

% ---------------------------------------------------------------------------
% 7.3 ComBat Harmonisation
% ---------------------------------------------------------------------------
\subsection{ComBat Harmonisation}
\label{sec:enc:combat}

Scanner-induced batch effects in multi-site MRI data can introduce spurious
variance in the latent space that confounds temporal growth patterns. We
assess and optionally apply ComBat harmonisation
\cite{johnson2007adjusting,fortin2018harmonization} to the encoded latent
vectors.

\subsubsection{Assessment Protocol}

Before applying ComBat, we assess the degree of scanner-induced shift using:
\begin{enumerate}[leftmargin=*]
  \item \textbf{MMD between cohorts:} $\text{MMD}^2(\mathbf{z}_{\text{BraTS}},
    \mathbf{z}_{\text{Andalusian}})$ with permutation test.
  \item \textbf{UMAP visualisation:} Coloured by site/scanner, to visually
    assess clustering by acquisition rather than biology.
  \item \textbf{Kolmogorov--Smirnov tests:} Per-dimension KS tests between
    sites.
\end{enumerate}

If the MMD permutation test $p$-value $< 0.05$, ComBat harmonisation is
applied. Otherwise, harmonisation is skipped to avoid unnecessary smoothing
of biological signal.

\subsubsection{ComBat Formulation}

ComBat models each latent dimension $j$ at site $s$ as:
\begin{equation}
  z_{ij} = \alpha_j + \mathbf{x}_i^\top \boldsymbol{\beta}_j
    + \gamma_{sj} + \delta_{sj}\,\epsilon_{ij},
  \label{eq:combat}
\end{equation}
where $\alpha_j$ is the overall mean, $\mathbf{x}_i$ are biological
covariates (age, sex if available), $\gamma_{sj}$ and $\delta_{sj}$ are
site-specific location and scale effects, and $\epsilon_{ij}$ is the
residual. Empirical Bayes estimation of $\gamma$ and $\delta$ provides
stable harmonisation even with small site sample sizes. The harmonised
value is:
\begin{equation}
  z_{ij}^* = \frac{z_{ij} - \hat{\alpha}_j - \mathbf{x}_i^\top
    \hat{\boldsymbol{\beta}}_j - \hat{\gamma}_{sj}}
    {\hat{\delta}_{sj}} + \hat{\alpha}_j
    + \mathbf{x}_i^\top\hat{\boldsymbol{\beta}}_j.
  \label{eq:combat_corrected}
\end{equation}

\subsubsection{Scanner Consistency Verification}

A critical prerequisite is verifying scanner consistency within each patient's
longitudinal series. If a patient was scanned on different scanners across
visits, standard ComBat (which assumes one site label per observation) may
distort temporal trends. In this case, LongComBat \cite{beer2020longcombat},
which explicitly models within-subject scanner changes in a mixed-effects
framework, should be used instead.

\begin{remark}
  ComBat's temporal preservation property requires careful verification.
  For a patient with timepoints $t_1, t_2$ scanned on the same scanner,
  the temporal difference $\Delta\mathbf{z} = \mathbf{z}(t_2) - \mathbf{z}(t_1)$
  should be preserved up to the site-specific scale factor $\hat{\delta}_s$.
  This is tested explicitly in the verification protocol (\cref{sec:evaluation}).
\end{remark}

% ---------------------------------------------------------------------------
% 7.4 Trajectory Construction
% ---------------------------------------------------------------------------
\subsection{Trajectory Construction}
\label{sec:enc:trajectories}

For each patient $i$ in the longitudinal cohort, we construct a temporal
trajectory by:
\begin{enumerate}[leftmargin=*]
  \item Encoding all $n_i$ MRI studies through the frozen pipeline, yielding
    $\{\mathbf{z}_i(t_{ij})\}_{j=1}^{n_i}$.
  \item Computing time offsets $t_{ij}$ in months from the patient's first
    scan: $t_{ij} = (\text{date}_{ij} - \text{date}_{i1}) / 30.44$.
  \item Verifying temporal ordering: $t_{i1} < t_{i2} < \cdots < t_{in_i}$.
  \item Extracting the partition subvectors:
    $\Vol(t_{ij}) = \mathbf{z}_i(t_{ij})[0:24]$, etc.
\end{enumerate}

The output is a set of per-patient trajectories:
\begin{equation}
  \mathcal{T}_i = \bigl\{
    \bigl(\mathbf{z}_i(t_{ij}),\; t_{ij}\bigr)
  \bigr\}_{j=1}^{n_i}, \quad i = 1, \ldots, N.
  \label{eq:trajectory}
\end{equation}

% ---------------------------------------------------------------------------
% 7.5 Cohort Distribution Comparison
% ---------------------------------------------------------------------------
\subsection{Cohort Distribution Comparison}
\label{sec:enc:distribution}

\begin{figure}[H]
  \centering
  \fbox{\parbox{0.8\textwidth}{\centering\vspace{3cm}
    \textbf{Figure placeholder:} Cohort distribution comparison.\\
    UMAP of BraTS-MEN (blue) vs.\ Andalusian (red) latent vectors.\\
    Left: before ComBat. Right: after ComBat (if applied).
  \vspace{3cm}}}
  \caption{UMAP comparison of BraTS-MEN ($n=1{,}000$) and Andalusian
    ($n=137$) latent representations. The degree of overlap indicates
    the compatibility of the two cohorts in the learned latent space.}
  \label{fig:cohort_umap}
\end{figure}

\begin{figure}[H]
  \centering
  \fbox{\parbox{0.8\textwidth}{\centering\vspace{3cm}
    \textbf{Figure placeholder:} Patient trajectories in latent space.\\
    2D UMAP with temporal arrows for 5--10 patients with $n_i \geq 3$.
  \vspace{3cm}}}
  \caption{Latent space trajectories for selected patients with $\geq 3$
    timepoints. Arrows indicate temporal progression. Smooth, directed
    trajectories suggest that the latent space captures meaningful temporal
    dynamics.}
  \label{fig:trajectories_umap}
\end{figure}


% --- Section 8: Phase 4 — GP-Based Growth Prediction ---
% =============================================================================
% Section 8: Growth Prediction — A Three-Stage Complexity Ladder
% =============================================================================
\section{Growth Prediction: A Three-Stage Complexity Ladder}
\label{sec:growth}

% ---------------------------------------------------------------------------
% 8.1 Overview and Design Rationale
% ---------------------------------------------------------------------------
\subsection{Overview and Design Rationale}
\label{sec:growth:overview}

This section presents the core methodological contribution: a three-stage
complexity ladder for meningioma growth prediction, where each stage must
justify its existence through improved out-of-sample performance over the
preceding one. The key design principle is parsimony: at $N = 31$--$58$
patients with $\sim$112 total observations, models with more than 2--3
predictor parameters per patient risk overfitting, and complex architectures
must demonstrate they capture signal rather than noise.

\subsubsection{Why Not a Neural ODE?}

The original pipeline design used a Neural ODE
\cite{chen2018neural,rubanova2019latent} for temporal modelling. This was
abandoned due to catastrophic overparameterisation: with 33~patients
(100~observations, 112~forward pairs), the Neural ODE's $\sim$3,100+
parameters produce a parameter-to-observation ratio of approximately 27.7.
The full pivot rationale is documented in \cite{pascual2026pivot}.

\subsubsection{Stage Summary}

The three stages, in order of increasing complexity:
\begin{description}[leftmargin=*,labelwidth=!]
  \item[Stage 1:] Segment tumour $\to$ extract volume $\to$ model scalar
    volume trajectories. Parameters: 2--5 (population level).
  \item[Stage 2:] Estimate patient-specific latent severity $s_i \in [0, 1]$
    $\to$ model growth as $g(s_i, t; \boldsymbol{\theta})$. Parameters:
    population $\boldsymbol{\theta}$ + 1 per patient.
  \item[Stage 3:] BrainSegFounder features $\to$ PCA $\to$ GP with ARD
    kernel. Parameters: population kernel hyperparameters + 3--6 PCA
    components.
\end{description}

A cross-cutting variance decomposition (\cref{sec:eval:variance}) quantifies
the marginal $\Delta R^2$ of each stage.

% ---------------------------------------------------------------------------
% 8.2 Stage 1: Volumetric Growth Models
% ---------------------------------------------------------------------------
\subsection{Stage 1: Segmentation-Based Volumetric Baseline}
\label{sec:growth:stage1}

Stage~1 is the simplest credible growth prediction approach: segment the
tumour from each MRI, compute whole-tumour volume in mm$^3$, and model the
resulting scalar volume trajectories. Three models of increasing flexibility
are evaluated.

\subsubsection{Model 1A: Scalar GP (ScalarGP)}

A single GP is fitted to the pooled volume observations across all patients:
\begin{equation}
  v_i(t_{ij}) \sim \GP\bigl(m(t),\; k_{5/2}(t, t') + \sigma_n^2 \delta(t, t')\bigr),
  \label{eq:scalargp}
\end{equation}
where $v_i(t_{ij}) = \log(V_{i,\text{WT}}(t_{ij}) + 1)$ is the
log-transformed whole-tumour volume, $m(t) = \beta_0 + \beta_1 t$ is a
linear mean function, and $k_{5/2}$ is the Mat\'{e}rn-5/2 kernel
(\cref{eq:matern52}).

\textbf{Hyperparameters:} $\{\beta_0, \beta_1, \sigma_f, \ell, \sigma_n\}$
= 5 total, estimated by maximising the log-marginal likelihood.

\textbf{Covariates:} Age and sex can be incorporated into the mean function:
\begin{equation}
  m_i(t) = \beta_0 + \beta_1 t + \beta_2 \,\text{age}_i
    + \beta_3 \,\text{sex}_i,
  \label{eq:scalargp_covariates}
\end{equation}
adding 2~parameters. Given the sample size constraints ($N/15 \approx 2$--$3$
allowed parameters), covariates are included only if they improve LOPO-CV
performance.

\subsubsection{Model 1B: Linear Mixed-Effects Model (LME)}

The LME introduces patient-specific random effects to separate population
and individual trends:
\begin{equation}
  v_i(t_{ij}) = \underbrace{(\beta_0 + b_{0i})}_{\text{intercept}}
    + \underbrace{(\beta_1 + b_{1i})}_{\text{slope}} \cdot t_{ij}
    + \epsilon_{ij},
  \label{eq:lme_vol}
\end{equation}
where $(\beta_0, \beta_1)$ are fixed effects (population intercept and
slope), $(b_{0i}, b_{1i}) \sim \mathcal{N}(\mathbf{0}, \boldsymbol{\Omega})$
are patient-specific random effects, and
$\epsilon_{ij} \sim \mathcal{N}(0, \sigma^2)$ is residual noise.

\textbf{Parameter count:} 2~fixed effects $(\beta_0, \beta_1)$, 3~variance
parameters (2~diagonal + 1~off-diagonal of $\boldsymbol{\Omega}$),
1~residual variance $(\sigma^2)$ = 6 total. Estimated via REML using
\texttt{statsmodels.MixedLM}.

\textbf{Single-patient inference:} The Best Linear Unbiased Predictor (BLUP)
\cite{robinson1991blup} provides patient-specific predictions with automatic
shrinkage: patients with few observations ($n_i = 2$) receive predictions
closer to the population mean, while patients with more observations
($n_i \geq 4$) receive data-driven predictions.

\subsubsection{Model 1C: Hierarchical GP (HGP)}

The HGP extends the LME by adding nonlinear residuals via a GP:
\begin{equation}
  v_i(t) \sim \GP\bigl(m(t),\; k_{5/2}(t, t') + \sigma_n^2 \delta(t, t')\bigr),
  \label{eq:hgp_vol}
\end{equation}
where the mean function is the LME population trend:
\begin{equation}
  m(t) = \hat{\beta}_0 + \hat{\beta}_1 \cdot t,
  \label{eq:gp_mean_lme}
\end{equation}
using fixed effects estimated from the LME. The GP captures nonlinear
deviations from the population trend.

\textbf{Hierarchical hyperparameter sharing:} Shared kernel hyperparameters
$\boldsymbol{\theta} = (\sigma_f, \ell, \sigma_n)$ are estimated by
maximising the pooled log-marginal likelihood across all patients:
\begin{equation}
  \hat{\boldsymbol{\theta}} = \arg\max_{\boldsymbol{\theta}}
    \sum_{i=1}^{N} \log p\bigl(v_i \mid
    \mathbf{t}_i, \boldsymbol{\theta}\bigr).
  \label{eq:pooled_marglik}
\end{equation}

\textbf{Parameter count:} 3~kernel hyperparameters (shared) + 2~mean
function parameters (from LME) = 5. Optimised via L-BFGS-B with 5~random
restarts.

\textbf{Model nesting:} The GP with linear mean function $\supset$ LME
(\cref{sec:bg:gp}), ensuring the HGP can never perform worse than the LME
on the training data. Improvement must come from the nonlinear Mat\'{e}rn
component, which we evaluate via log-marginal likelihood comparison.

\subsubsection{Gompertz Mean Function}

As an alternative to the linear population mean, we evaluate a Gompertz
mean function motivated by the growth biology literature
(\cref{sec:bg:growth_biology}):
\begin{equation}
  m_G(t) = V_0 \exp\!\left(\frac{\alpha}{\beta}
    \bigl(1 - e^{-\beta t}\bigr)\right),
  \label{eq:gompertz_mean}
\end{equation}
where $V_0$ is the initial volume, $\alpha$ is the initial growth rate,
and $\beta$ is the deceleration parameter. In log-volume space this becomes
$\log m_G(t) = \log V_0 + \frac{\alpha}{\beta}(1 - e^{-\beta t})$. The
Gompertz mean adds 1~parameter ($\beta$) beyond the linear mean, and is
evaluated as ablation~A2 (\cref{sec:eval:ablations}).

% ---------------------------------------------------------------------------
% 8.3 Stage 2: Latent Severity Model
% ---------------------------------------------------------------------------
\subsection{Stage 2: Latent Severity Model}
\label{sec:growth:stage2}

Stage~2 posits that inter-patient growth heterogeneity is largely explained
by a single latent severity variable $s_i \in [0, 1]$. This is motivated by
Vaghi et al.'s \cite{vaghi2020reduced} finding that a single patient-specific
parameter captures most growth variability in tumour populations.

\subsubsection{Quantile Transform}

To place all patients on a common scale, we apply a quantile transform to
both time and growth:
\begin{enumerate}[leftmargin=*]
  \item \textbf{Temporal normalisation.} Within each patient, times are
    normalised to $[0, 1]$:
    $\tilde{t}_{ij} = (t_{ij} - t_{i,\min}) / (t_{i,\max} - t_{i,\min})$.
  \item \textbf{Growth normalisation.} Volume changes are mapped to
    quantiles of the population distribution:
    $q_{ij} = F_{\text{pop}}(\Delta v_{ij})$, where $F_{\text{pop}}$ is the
    empirical CDF of all observed volume changes across the training set.
\end{enumerate}

The quantile-transformed observations $q_{ij} \in [0, 1]$ represent the
relative position of each observation within the population growth
distribution.

\subsubsection{NLME Formulation}

The severity model is a nonlinear mixed-effects model:
\begin{equation}
  q_{ij} = g(s_i, \tilde{t}_{ij}; \boldsymbol{\theta}) + \epsilon_{ij},
  \quad \epsilon_{ij} \sim \mathcal{N}(0, \sigma^2),
  \label{eq:severity_nlme}
\end{equation}
where $g$ is a monotonic growth function satisfying boundary conditions
$g(s, 0) = 0$ (no growth at baseline) and $g(s, 1) \leq s$ (severity bounds
maximum growth). The latent severity $s_i \sim \text{Beta}(\alpha, \beta)$ has
a population prior.

\subsubsection{Growth Function Options}

Three growth functions are evaluated for $g$:

\begin{description}[leftmargin=*,labelwidth=!]
  \item[Sigmoid.] The simplest parametric option:
    $g(s, t; \boldsymbol{\theta}) = s \cdot \sigma(\theta_1 (t - \theta_2))$,
    where $\sigma$ is the logistic sigmoid and
    $\boldsymbol{\theta} = (\theta_1, \theta_2)$ controls the steepness and
    inflection point. Monotonic in $s$ by construction. 2~population
    parameters.

  \item[Reduced Gompertz.] Following Vaghi et al.\
    \cite{vaghi2020reduced}:
    $g(s, t; \beta) = s \cdot [1 - \exp(-\alpha(s)(1 - e^{-\beta t})/\beta)]$,
    where $\alpha(s)$ maps severity to growth rate. 1~population parameter
    ($\beta$).

  \item[CMNN.] A constrained monotonic neural network
    \cite{runje2023cmnn} with non-negative weights ensuring
    $\partial g / \partial s \geq 0$ everywhere. The most flexible option
    but with the highest parameter count ($\sim$10--20 weights). Evaluated
    only if the parametric options prove insufficient.
\end{description}

These are compared in ablation~A3 (\cref{sec:eval:ablations}).

\subsubsection{Joint Estimation}

The population parameters $\boldsymbol{\theta}$ and patient severities
$\{s_i\}_{i=1}^N$ are estimated jointly by maximising the marginal
likelihood:
\begin{equation}
  \hat{\boldsymbol{\theta}} = \arg\max_{\boldsymbol{\theta}}
    \sum_{i=1}^{N} \log \int p(\mathbf{q}_i \mid s_i, \boldsymbol{\theta})
    \, p(s_i) \, ds_i.
  \label{eq:severity_mle}
\end{equation}
For the Beta prior, the inner integral is computed via Gauss--Hermite
quadrature or Laplace approximation.

\subsubsection{Test-Time Severity Estimation}

For a held-out patient with observed trajectory
$\{(t_{ij}, q_{ij})\}_{j=1}^{n_i}$, the posterior severity is:
\begin{equation}
  p(s_i \mid \mathbf{q}_i, \hat{\boldsymbol{\theta}}) \propto
    p(\mathbf{q}_i \mid s_i, \hat{\boldsymbol{\theta}}) \, p(s_i),
  \label{eq:severity_posterior}
\end{equation}
which is maximised or integrated over for point and distributional predictions
respectively. For patients with $n_i = 2$ (the most common case), the
posterior is dominated by the prior, yielding predictions close to the
population median---a desirable behaviour analogous to LME shrinkage.

\subsubsection{Severity Estimation from Baseline Features}

A secondary objective of Stage~2 is to predict severity from baseline
(single-timepoint) information, enabling predictions without longitudinal
data. Candidate predictors include:
\begin{itemize}[leftmargin=*]
  \item Clinical covariates: age, sex, tumour location.
  \item Baseline volume $V_{i,0}$.
  \item BrainSegFounder features at $t_0$ (linking to Stage~3).
\end{itemize}
This is evaluated in ablation~A4 (\cref{sec:eval:ablations}).

% ---------------------------------------------------------------------------
% 8.4 Stage 3: Representation-Learning Augmented Prediction
% ---------------------------------------------------------------------------
\subsection{Stage 3: Representation-Learning Augmented Prediction}
\label{sec:growth:stage3}

Stage~3 asks whether high-dimensional deep features from BrainSegFounder
capture growth-relevant information beyond what volume and severity already
explain. This stage uses the LoRA-adapted encoder (\cref{sec:lora}) and
SDP (\cref{sec:sdp}) from the original pipeline, but positions them as the
final complexity level rather than the primary approach.

\subsubsection{Feature Extraction and Compression}

The frozen pipeline (encoder + SDP) produces per-scan latent vectors
$\mathbf{z} \in \R^{128}$. For growth prediction, these are compressed via
within-fold PCA to $\tilde{\mathbf{z}} \in \R^{d_{\text{PCA}}}$ with
$d_{\text{PCA}} \in \{3, 4, 5, 6\}$ selected by the sample size constraint
and cross-validated performance. PCA is fitted on the training set within
each LOPO-CV fold to prevent information leakage.

\subsubsection{GP with ARD Kernel}

The compressed features are used as inputs to a GP with an ARD kernel:
\begin{equation}
  k_{\text{ARD}}(\tilde{\mathbf{z}}, \tilde{\mathbf{z}}'; t, t') =
    k_{5/2}(t, t') \cdot \prod_{d=1}^{d_{\text{PCA}}}
    k_{\text{SE}}\!\left(\frac{\tilde{z}^{(d)} - \tilde{z}'^{(d)}}{\ell_d}\right)
    + \sigma_n^2 \delta_{ii'} \delta(t, t'),
  \label{eq:ard_kernel}
\end{equation}
where $\ell_d$ are per-dimension length-scales. Dimensions with large
$\ell_d$ are effectively deactivated, providing automatic feature selection.
The mean function is the Stage~1 population trend.

\textbf{Parameter count:} $d_{\text{PCA}}$ ARD length-scales + 3 temporal
kernel hyperparameters + mean function parameters. For $d_{\text{PCA}} = 4$:
$\sim$9 total hyperparameters, estimated from the pooled marginal likelihood.

\subsubsection{Integration with Earlier Stages}

Stage~3 does not replace Stages~1 and~2 but augments them. The deep features
enter as additional covariates in the GP mean function or as ARD inputs
alongside time, while the population trend from Stage~1 remains the prior
mean. The marginal contribution is quantified by the variance decomposition
(\cref{sec:eval:variance}).

\subsubsection{SDP with Supervised Residual Targets}

An alternative to PCA compression is to use the SDP's supervised partition
$\Vol \in \R^{32}$ directly. The SDP is designed to place volume-relevant
information in a structured subspace. However, 32~dimensions is still far
too many for GP modelling at $N = 33$. Two options are evaluated:
\begin{enumerate}[leftmargin=*]
  \item PCA on $\Vol$ alone (expected to require fewer components than PCA
    on full $\mathbf{z}$, since the SDP has already concentrated
    volume-relevant information).
  \item Using $\Vol$ as input to a GP with an ICM coregionalisation matrix
    (\cref{eq:icm}), treating each latent dimension as an output channel.
    This is the PA-MOGP structure from the original design.
\end{enumerate}

\begin{remark}[Stage~3 is conditional]
  Stage~3 is attempted only if Stages~1 and~2 leave unexplained variance
  that could plausibly be captured by additional features. If the variance
  decomposition shows that volume and severity together explain $R^2 > 0.70$,
  the marginal room for improvement from deep features is limited, and
  Stage~3 serves primarily as a negative result confirming parsimony.
\end{remark}

% ---------------------------------------------------------------------------
% 8.5 LOPO-CV Protocol
% ---------------------------------------------------------------------------
\subsection{Leave-One-Patient-Out Cross-Validation}
\label{sec:growth:lopo}

All three stages share a common LOPO-CV protocol. With only 31--58~patients,
$k$-fold CV leaves too few patients per fold for stable model estimation.
LOPO-CV uses maximum training data per fold:

\begin{algorithm}[H]
  \caption{LOPO-CV Evaluation Protocol (shared across all stages)}
  \label{alg:lopo}
  \begin{algorithmic}[1]
    \FOR{$i = 1$ to $N$}
      \STATE $\mathcal{D}_{\text{train}} \leftarrow
        \{\mathcal{T}_j : j \neq i\}$ \COMMENT{$N-1$ patients}
      \STATE $\mathcal{D}_{\text{test}} \leftarrow \{\mathcal{T}_i\}$
        \COMMENT{held-out patient}
      \STATE \textbf{Within-fold operations:}
      \STATE \quad Fit PCA on $\mathcal{D}_{\text{train}}$ (Stage~3 only)
      \STATE \quad Estimate severity $\{s_j\}_{j \neq i}$ (Stage~2 only)
      \STATE \quad Fit model on $\mathcal{D}_{\text{train}}$
      \STATE \quad Predict at patient $i$'s observation times
      \STATE \quad Record per-patient metrics
    \ENDFOR
    \STATE Aggregate metrics across $N$ folds
    \STATE Compute .632+ bootstrap CIs
  \end{algorithmic}
\end{algorithm}

All model fitting, feature transformation (PCA), and severity estimation
are performed within each fold. This ensures that LOPO-CV metrics reflect
true out-of-sample performance without information leakage.

For each held-out patient $i$, predictions are evaluated at the patient's
observed timepoints. For Stage~2, the held-out patient's severity is
re-estimated from their observations using the population model fitted on the
remaining patients (\cref{eq:severity_posterior}).

% ---------------------------------------------------------------------------
% 8.6 Volume Decoding and Uncertainty
% ---------------------------------------------------------------------------
\subsection{Volume Representation and Uncertainty}
\label{sec:growth:decoder}

\subsubsection{Volume Representation}

All models in Stages~1 and~2 operate on log-transformed whole-tumour volume:
\begin{equation}
  v = \log(V_{\text{WT}} + 1),
  \label{eq:log_vol}
\end{equation}
where $V_{\text{WT}}$ is the whole-tumour volume in mm$^3$. This is the
single clinically relevant growth endpoint for meningiomas. The
log-transformation stabilises variance and makes growth rates approximately
additive.

For Stage~3, the SDP semantic head decodes latent predictions to
log-volumes:
\begin{equation}
  \hat{v} = \pi_{\text{vol}}(\hat{\Vol})
    = \mathbf{W}_v \hat{\Vol} + \mathbf{b}_v,
  \quad \hat{v} \in \R^1,
  \label{eq:vol_decode}
\end{equation}
with de-normalisation $V_{\text{pred}} = \exp(\hat{v} \cdot \sigma_{\text{vol}}
+ \mu_{\text{vol}}) - 1$.

\subsubsection{Uncertainty Propagation}

For GP models that provide predictive variance $\sigma^2_{\hat{v}}$ in
log-volume space, the physical volume uncertainty is propagated through the
exponential via the delta method:
\begin{equation}
  \sigma_{V_{\text{pred}}} \approx V_{\text{pred}} \cdot \sigma_{\hat{v}}
    \cdot \sigma_{\text{vol}},
  \label{eq:delta_method}
\end{equation}
where the multiplicative structure reflects the log-normal nature of volume
uncertainty. For prediction intervals, we report the 95\% credible interval
from the GP posterior.

% ---------------------------------------------------------------------------
% 8.7 Failure Recovery
% ---------------------------------------------------------------------------
\subsection{Failure Recovery Protocol}
\label{sec:growth:failure}

If GP hyperparameter fitting diverges or produces degenerate results, the
following recovery steps are applied in order:
\begin{enumerate}[leftmargin=*]
  \item Add diagonal jitter to the kernel matrix:
    $\mathbf{K} \leftarrow \mathbf{K} + 10^{-6}\,\mathbf{I}$.
  \item Constrain length-scale bounds: $\ell \in [1, 120]$ months.
  \item Increase L-BFGS-B random restarts from 5 to 10.
  \item Fall back to a simpler kernel (Mat\'{e}rn-5/2 $\to$ SE).
  \item If all fail: report per-dimension diagnostics and halt.
\end{enumerate}

For Stage~2 severity estimation, if the joint optimisation fails to converge:
\begin{enumerate}[leftmargin=*]
  \item Reduce the growth function complexity (CMNN $\to$ sigmoid).
  \item Initialise severities from the empirical growth rate ranking.
  \item Fix population parameters and estimate severities only (profile
    likelihood).
\end{enumerate}


% --- Section 9: Evaluation ---
% =============================================================================
% Section 9: Evaluation
% =============================================================================
\section{Evaluation}
\label{sec:evaluation}

% ---------------------------------------------------------------------------
% 9.1 Quality Targets
% ---------------------------------------------------------------------------
\subsection{Quality Targets}
\label{sec:eval:quality}

The pipeline defines quantitative quality targets at each phase, with minimum
thresholds that block further progress and aspirational targets for optimal
performance.

\begin{table}[H]
  \centering
  \caption{Phase 1 quality targets (encoder adaptation).}
  \label{tab:quality_phase1}
  \begin{tabular}{lccc}
    \toprule
    \textbf{Metric} & \textbf{Target} & \textbf{Minimum} & \textbf{Status} \\
    \midrule
    Dice (WT) on \texttt{lora\_val} & $\geq 0.85$ & $\geq 0.80$ & \textit{[to fill]} \\
    Dice improvement over frozen BSF & $\geq 0.05$ & $\geq 0.02$ & \textit{[to fill]} \\
    Linear probe Vol $R^2$ (adapted) & $\geq 0.50$ & $\geq 0.30$ & \textit{[to fill]} \\
    \bottomrule
  \end{tabular}
\end{table}

\begin{table}[H]
  \centering
  \caption{Phase 2 quality targets (SDP).}
  \label{tab:quality_phase2}
  \begin{tabular}{lccc}
    \toprule
    \textbf{Metric} & \textbf{Target} & \textbf{Minimum} & \textbf{Status} \\
    \midrule
    Vol $R^2$ on validation & $\geq 0.90$ & $\geq 0.80$ & \textit{[to fill]} \\
    Loc $R^2$ on validation & $\geq 0.95$ & $\geq 0.85$ & \textit{[to fill]} \\
    Shape $R^2$ on validation & $\geq 0.40$ & $\geq 0.25$ & \textit{[to fill]} \\
    Max cross-partition corr & $< 0.20$ & $< 0.30$ & \textit{[to fill]} \\
    Dims with var $> 0.5$ & $\geq 95\%$ & $\geq 85\%$ & \textit{[to fill]} \\
    $\dCor(\text{vol}, \text{loc})$ & $< 0.10$ & $< 0.20$ & \textit{[to fill]} \\
    \bottomrule
  \end{tabular}
\end{table}

\begin{table}[H]
  \centering
  \caption{Phase 4 quality targets (growth prediction).}
  \label{tab:quality_phase4}
  \begin{tabular}{lccc}
    \toprule
    \textbf{Metric} & \textbf{Target} & \textbf{Minimum} & \textbf{Status} \\
    \midrule
    Vol prediction $R^2$ (LOPO, best model) & $\geq 0.70$ & $\geq 0.50$ & \textit{[to fill]} \\
    Calibration (95\% CI coverage, GP models) & 0.90--0.98 & $\geq 0.80$ & \textit{[to fill]} \\
    Per-patient trajectory $r$ ($n_i \geq 3$) & $\geq 0.80$ & $\geq 0.60$ & \textit{[to fill]} \\
    PA-MOGP $R^2$ $>$ H-GP $R^2$ & $> 0$ & --- & \textit{[to fill]} \\
    \bottomrule
  \end{tabular}
\end{table}

% ---------------------------------------------------------------------------
% 9.2 Ablation Matrix
% ---------------------------------------------------------------------------
\subsection{Ablation Study Design}
\label{sec:eval:ablations}

Ten ablation experiments systematically evaluate design decisions across all
pipeline phases:

\begin{table}[H]
  \centering
  \caption{Full ablation matrix (A1--A10).}
  \label{tab:ablation_matrix}
  \begin{tabular}{clcp{5.5cm}}
    \toprule
    \textbf{ID} & \textbf{Variable} & \textbf{Conditions} & \textbf{Primary metric} \\
    \midrule
    A1 & LoRA rank & $\{2, 4, 8, 16, 32\}$ & Dice, probe $R^2$ \\
    A2 & LoRA vs.\ DoRA & $\{\LoRA, \text{DoRA}\}$ at $r=8$ & Dice, probe $R^2$ \\
    A3 & Aux semantic heads & $\{$with, without$\}$ & Phase~2 $R^2$ \\
    A4 & SDP dimension & $\{64, 128, 256\}$ & $R^2$, dCor \\
    A5 & VICReg + dCor & $\{$full, no cov, no dCor, no both$\}$ & Cross-partition corr \\
    A6 & Growth model & $\{\LME, \HGP, \PAMOGP\}$ & Vol $R^2$ (LOPO) \\
    A7 & ComBat effect & $\{$with, without$\}$ & Vol $R^2$ (LOPO) \\
    A8 & GP mean function & $\{$zero, linear, Gompertz$\}$ & Vol $R^2$ (LOPO) \\
    A9 & GP kernel & $\{$M-3/2, M-5/2, SE$\}$ & Vol $R^2$ (LOPO) \\
    A10 & Cross-partition coupling & $\{$with, without$\}$ & Vol $R^2$, calibration \\
    \bottomrule
  \end{tabular}
\end{table}

\subsubsection{A6: Model Comparison}

The head-to-head comparison of the three GP models is the central evaluation
of the pipeline. The comparison uses LOPO-CV and reports:

\begin{description}[leftmargin=*]
  \item[Vol $R^2$:] Coefficient of determination between predicted and
    observed whole-tumour volume across all held-out observations.
  \item[Vol MAE:] Mean absolute error in physical volume (mm$^3$).
  \item[Latent MSE:] Mean squared error in the latent volume partition
    $\Vol$ (before decoding).
  \item[Calibration:] Fraction of true observations within the predicted
    95\% confidence interval (target: 0.90--0.98).
  \item[Per-patient $r$:] Pearson correlation between predicted and observed
    volume trajectories for patients with $n_i \geq 3$.
\end{description}

\begin{table}[H]
  \centering
  \caption{Model comparison results (LOPO-CV, 33 folds).}
  \label{tab:model_comparison}
  \begin{tabular}{lccccc}
    \toprule
    \textbf{Model} & \textbf{Params} & \textbf{Vol $R^2$} &
      \textbf{Vol MAE} & \textbf{Calibration} & \textbf{Per-patient $r$} \\
    \midrule
    LME     & 144 & \textit{[to fill]} & \textit{[to fill]} & --- & \textit{[to fill]} \\
    H-GP    & 72  & \textit{[to fill]} & \textit{[to fill]} & \textit{[to fill]} & \textit{[to fill]} \\
    PA-MOGP & 98  & \textit{[to fill]} & \textit{[to fill]} & \textit{[to fill]} & \textit{[to fill]} \\
    \bottomrule
  \end{tabular}
\end{table}

% ---------------------------------------------------------------------------
% 9.3 Figures
% ---------------------------------------------------------------------------
\subsection{Publication Figures}
\label{sec:eval:figures}

The evaluation produces 13 publication-quality figures documenting each
pipeline phase. Key figures include:

\begin{figure}[H]
  \centering
  \fbox{\parbox{0.8\textwidth}{\centering\vspace{3cm}
    \textbf{Figure placeholder:} Volume prediction scatter.\\
    Predicted vs.\ actual $\Delta V$ for all LOPO-CV test pairs,\\
    coloured by model (LME: blue, H-GP: orange, PA-MOGP: green).
  \vspace{3cm}}}
  \caption{Predicted vs.\ actual volume change ($\Delta V$, mm$^3$) across
    all LOPO-CV held-out observation pairs, coloured by model. The
    diagonal represents perfect prediction. Dispersion around the diagonal
    quantifies prediction error.}
  \label{fig:vol_scatter}
\end{figure}

\begin{figure}[H]
  \centering
  \fbox{\parbox{0.8\textwidth}{\centering\vspace{3cm}
    \textbf{Figure placeholder:} Trajectory prediction with uncertainty.\\
    3--4 example patients showing predicted trajectory\\
    (mean $\pm$ 95\% CI) overlaid on actual observations.
  \vspace{3cm}}}
  \caption{Growth trajectory predictions for selected patients. Solid lines
    show GP posterior mean; shaded regions show 95\% confidence intervals.
    Dots are actual observations. The GP uncertainty naturally widens with
    extrapolation distance.}
  \label{fig:trajectory_predictions}
\end{figure}

\begin{figure}[H]
  \centering
  \fbox{\parbox{0.8\textwidth}{\centering\vspace{3cm}
    \textbf{Figure placeholder:} Model comparison bar chart.\\
    $R^2$, MAE, and calibration for each model.
  \vspace{3cm}}}
  \caption{Head-to-head model comparison across three metrics. Error bars
    represent bootstrap 95\% confidence intervals computed over the 33
    LOPO-CV folds.}
  \label{fig:model_comparison}
\end{figure}

\begin{figure}[H]
  \centering
  \fbox{\parbox{0.8\textwidth}{\centering\vspace{3cm}
    \textbf{Figure placeholder:} Disentanglement matrix.\\
    Cross-partition correlation heatmap from SDP.
  \vspace{3cm}}}
  \caption{Cross-partition correlation heatmap after SDP training. Each cell
    shows the maximum absolute Pearson correlation between dimensions in the
    corresponding partition pair. Diagonal blocks (within-partition) show
    internal structure; off-diagonal blocks should approach zero for
    successful disentanglement.}
  \label{fig:disentanglement_matrix}
\end{figure}

\begin{figure}[H]
  \centering
  \fbox{\parbox{0.8\textwidth}{\centering\vspace{3cm}
    \textbf{Figure placeholder:} Learned GP length-scales per partition.\\
    Revealing characteristic timescales of volume/location/shape dynamics.
  \vspace{3cm}}}
  \caption{Learned GP kernel length-scales ($\ell$) for the PA-MOGP model.
    Volume dynamics have the longest characteristic timescale (slow growth),
    while shape dynamics are faster (shorter $\ell$). Location length-scales
    are expected to be very long (near-static centroid).}
  \label{fig:lengthscales}
\end{figure}

% ---------------------------------------------------------------------------
% 9.4 Methodology Revision (R1)
% ---------------------------------------------------------------------------
\subsection{R1 Methodology Revision}
\label{sec:eval:revision}

Empirical findings from Phases~1--2 motivated three principled revisions to
the pipeline architecture, collectively designated the R1 revision.

\subsubsection{Shape Partition Failure}

Across all experimental conditions, the shape partition achieved
$R^2 \leq 0.11$ for linear probes and approximately zero for RBF probes.
Cross-domain shape transfer was uniformly catastrophic ($R^2 = -2.00$).
This failure is biologically grounded: gliomas are infiltrative and
irregular, while meningiomas are well-circumscribed and dural-based. Shape
semantics differ qualitatively between the two tumour types, making a shared
shape representation a category error.

\begin{table}[H]
  \centering
  \caption{Shape probe $R^2$ across experimental conditions.}
  \label{tab:shape_failure}
  \begin{tabular}{lccl}
    \toprule
    \textbf{Condition} & \textbf{Linear $R^2$} & \textbf{RBF $R^2$} &
      \textbf{Cross-domain} \\
    \midrule
    Frozen baseline   & 0.08  & $\approx 0$ & $-2.00$ \\
    MEN LoRA $r$=8    & 0.11  & $\approx 0$ & $-2.00$ \\
    Dual LoRA $r$=8   & $-0.06$ & $\approx 0.15$ & $-2.00$ \\
    \bottomrule
  \end{tabular}
\end{table}

\textbf{Decision:} Drop the shape partition entirely from the SDP.

\subsubsection{Location: Temporally Static}

Meningiomas are extra-axial, dural-attached tumours that do not migrate. The
centroid coordinates at $t_0$ are essentially identical at subsequent
timepoints (modulo asymmetric growth). Modelling location as a time-varying
GP dimension fits a constant function, which is statistically vacuous and
wastes degrees of freedom. Within-domain location $R^2$ is reasonable
(0.35 MEN, 0.68 GLI), confirming the information is present but does not
need temporal modelling.

\textbf{Decision:} Relocate location from temporal latent dimension to static
GP covariate, entering the mean function as
$m_i(t) = \beta_0 + \beta_1 t + \boldsymbol{\gamma}^\top \mathbf{c}_i$,
where $\mathbf{c}_i$ is the centroid from patient $i$'s initial scan.

\subsubsection{Volume Target Simplification}

The 4-dimensional volume target $[\log V_{\text{total}}, \log V_{\text{NCR}},
\log V_{\text{ED}}, \log V_{\text{ET}}]$ is problematic because: (i)~label
semantics differ across tumour types (GLI oedema $\neq$ MEN oedema);
(ii)~sub-compartment segmentation quality is poor (TC Dice 0.40 on MEN);
(iii)~whole-tumour volume is the only clinically relevant growth endpoint for
meningiomas.

\textbf{Decision:} Reduce to a single target
$v^* = \log(V_{\text{WT}} + 1) \in \R^1$.

\subsubsection{Revised Architecture Summary}

\begin{equation}
  \text{R1:}\quad \mathbf{z} \in \R^{128} =
    [\underbrace{\Vol \in \R^{32}}_{\text{supervised}}
    \mid \underbrace{\Res \in \R^{96}}_{\text{unsupervised}}],
  \quad \pi_{\text{vol}}: \R^{32} \to \R^1.
  \label{eq:r1_partition}
\end{equation}

The PA-MOGP simplifies to a standard MOGP on $\R^{32}$ with a single
Mat\'{e}rn-5/2 kernel and ICM coregionalisation matrix
$\mathbf{B} = \mathbf{W}\mathbf{W}^\top + \diag(\boldsymbol{\kappa})$.

% ---------------------------------------------------------------------------
% 9.5 LoRA Marginality Discussion
% ---------------------------------------------------------------------------
\subsection{LoRA Marginality}
\label{sec:eval:marginality}

A notable finding is the marginal benefit of LoRA adaptation for downstream
feature quality. The frozen BrainSegFounder baseline achieves
$R^2 = 0.794$ for volume probing, while the best LoRA condition adds only
$\Delta R^2 = +0.016$ (within noise). This suggests that BrainSegFounder's
SSL pretraining on 41,400+ subjects already produces features with strong
meningioma volume decodability.

Several hypotheses explain this observation:
\begin{enumerate}[leftmargin=*]
  \item \textbf{SSL pretraining captures domain-general features.} The UK
    Biobank pretraining stage exposes the encoder to massive anatomical
    diversity, producing features that generalise across brain pathologies.
  \item \textbf{LoRA's low-rank constraint limits adaptation.} With
    $r = 8$ and $<0.32\%$ trainable parameters, LoRA may not have
    sufficient capacity to substantially restructure the 768-dimensional
    feature space.
  \item \textbf{Segmentation is a weak proxy for semantic quality.}
    Improving segmentation Dice does not necessarily improve the linear
    decodability of aggregate semantic properties (volume, location).
\end{enumerate}

This finding has important practical implications: for meningioma growth
prediction, the frozen BrainSegFounder encoder may be a sufficient starting
point, with the SDP and GP models providing the primary value.


% --- Section 10: Conclusion ---
% =============================================================================
% Section 10: Conclusion
% =============================================================================
\section{Conclusion}
\label{sec:conclusion}

% ---------------------------------------------------------------------------
% 10.1 Summary of Contributions
% ---------------------------------------------------------------------------
\subsection{Summary of Contributions}
\label{sec:conc:summary}

This thesis has presented a three-stage complexity ladder for meningioma
growth forecasting from longitudinal multi-modal 3D MRI. Rather than
proposing a single complex pipeline and reporting its metrics, we
systematically evaluate whether each successive level of model
sophistication provides measurable out-of-sample improvement at $N = 31$--$58$
patients. The key contributions and findings are:

\begin{enumerate}[leftmargin=*]
  \item \textbf{The first out-of-sample meningioma growth prediction
    benchmark.} Using LOPO-CV on a longitudinal cohort, we provide the
    first reported $R^2$, MAE, and calibration metrics for meningioma
    volume prediction. The volumetric baseline (Stage~1: ScalarGP, LME,
    HGP) establishes a reference point that, to our knowledge, did not
    previously exist in the literature. The HGP with population linear mean
    from the LME provides calibrated uncertainty with automatic shrinkage
    toward the population trend for patients with sparse observations.

  \item \textbf{A latent severity framework for growth characterisation.}
    Stage~2 formulates meningioma growth heterogeneity as a latent variable
    problem connected to Item Response Theory and the reduced Gompertz model
    \cite{vaghi2020reduced}. A single patient-specific severity parameter
    $s_i \in [0, 1]$ captures inter-patient variability while respecting
    the sample size constraint of 2--3 parameters per predictor. The
    framework provides interpretable patient stratification aligned with
    the three growth subgroups identified by Engelhardt et al.\
    \cite{engelhardt2023gompertz}.

  \item \textbf{Systematic assessment of representation learning value.}
    By placing the full BrainSegFounder + LoRA + SDP pipeline as Stage~3,
    we directly quantify whether 768-dimensional deep features (compressed
    via PCA to 3--6 dimensions) capture growth-relevant information beyond
    what volume and severity already explain. The marginal benefit of LoRA
    adaptation for feature quality ($\Delta R^2 = +0.016$ on semantic
    probes) motivates scepticism that the full representation learning
    pipeline is necessary for this task.

  \item \textbf{Variance decomposition as an analytical framework.}
    The cross-cutting analysis decomposes LOPO-CV prediction variance across
    the three stages using paired permutation tests and BCa bootstrap CIs.
    This provides not only a model selection criterion but also a
    methodological template for evaluating complexity--benefit tradeoffs in
    small-sample longitudinal prediction problems more broadly.

  \item \textbf{Empirically-driven methodology revision.} Rigorous
    evaluation revealed that shape disentanglement fails ($R^2 \leq 0.11$),
    meningioma location is temporally static, and sub-compartment volume
    targets are label-dependent and clinically irrelevant. These findings
    motivated a principled simplification and the reframing from a linear
    pipeline to the complexity ladder.
\end{enumerate}

% ---------------------------------------------------------------------------
% 10.2 Limitations
% ---------------------------------------------------------------------------
\subsection{Limitations}
\label{sec:conc:limitations}

\subsubsection{Cohort Size}

The longitudinal cohort ($N = 31$--$58$ patients, $\sim$112~observations)
is the primary bottleneck. With approximately 57.6\% of patients having
only 2~studies, per-patient growth trajectories are sparse. The GP and LME
frameworks handle this via automatic shrinkage toward the population mean,
but predictions for 2-observation patients are necessarily dominated by
population-level dynamics rather than individual-specific growth rates.
The sample size also limits the statistical power of the variance
decomposition: Varoquaux \cite{varoquaux2018cv} showed that CV error
estimates have standard deviations of $\sim$$\pm$10\% at $N = 100$, and
uncertainty is larger still at $N = 33$.

\subsubsection{Ordinal Time}

At the time of writing, exact acquisition timestamps for the longitudinal
cohort are pending. If only ordinal timepoint indices are available (scan~1,
scan~2, scan~3, \ldots) rather than calendar dates, the temporal models
operate on a surrogate time axis. This does not invalidate the complexity
ladder comparison (all models use the same time representation) but limits
the clinical interpretability of growth rate estimates and length-scale
hyperparameters.

\subsubsection{LoRA Marginality and Stage~3 Viability}

The marginal benefit of LoRA adaptation ($\Delta R^2 = +0.016$ for
semantic probing) suggests that the entire Stage~3 pipeline may provide
limited value at the current sample size. This is, however, a finding
rather than a failure: establishing that simple volumetric models suffice
is a valuable negative result that can save computational resources in
clinical deployment.

\subsubsection{Single-Site Longitudinal Data}

The longitudinal cohort is collected from hospitals in a single region
(Andalusia, Spain), which may limit generalisability to other populations,
scanner manufacturers, and acquisition protocols. ComBat harmonisation
addresses scanner batch effects but cannot correct for biological population
differences.

\subsubsection{Absence of Clinical Validation}

The pipeline has not been validated against clinical outcomes (intervention
decisions, actual tumour behaviour at longer follow-up). Volume prediction
$R^2$ and calibration metrics are necessary but not sufficient for clinical
utility. A positive growth prediction does not, by itself, determine the
appropriate clinical action.

% ---------------------------------------------------------------------------
% 10.3 Future Work
% ---------------------------------------------------------------------------
\subsection{Future Work}
\label{sec:conc:future}

\subsubsection{Cohort Expansion to $N \geq 50$}

The most impactful improvement would be expanding the longitudinal cohort.
Multi-centre collaborations targeting $N \geq 50$ patients would: (i)~tighten
the variance decomposition CIs, enabling confident $\Delta R^2$ conclusions;
(ii)~potentially unlock value in Stage~2 and~3 by providing sufficient data
for more complex models; (iii)~enable subgroup analyses by tumour grade,
location, and patient demographics.

\subsubsection{Exact Timestamps}

Obtaining exact acquisition dates would enable physically meaningful
growth rate estimation, improve GP length-scale interpretability, and allow
direct comparison with the volumetric growth rates reported in the clinical
literature \cite{hashimoto2012natural,nakamura2003natural,
engelhardt2023gompertz}.

\subsubsection{Clinical Decision Support Integration}

The GP posterior's calibrated uncertainty is directly suitable for clinical
decision support: the probability that a tumour will exceed a volume
threshold $V_{\text{critical}}$ within a specified time horizon can be
computed analytically from the Gaussian predictive distribution:
\begin{equation}
  P(V(t_{\text{future}}) > V_{\text{critical}}) =
    1 - \Phi\!\left(\frac{\log(V_{\text{critical}} + 1) - \mu_*(t_{\text{future}})}
    {\sigma_*(t_{\text{future}})}\right),
  \label{eq:clinical_decision}
\end{equation}
where $\Phi$ is the standard normal CDF and $(\mu_*, \sigma_*)$ are the GP
posterior mean and standard deviation in log-volume space. Integrating this
into a clinical workflow would require prospective validation and regulatory
approval.

\subsubsection{Severity-Guided Surveillance}

If the severity model (Stage~2) provides robust patient stratification,
it could inform surveillance scheduling: high-severity patients receive
shorter follow-up intervals, while low-severity patients can be monitored
less frequently. This requires validation of the severity--growth
relationship at longer follow-up horizons.

\subsubsection{Foundation Model Scaling}

BrainSegFounder is a relatively small foundation model (BSF-Tiny, 62M
parameters). Larger variants or more recent foundation models pretrained on
broader datasets may provide stronger initial features, potentially shifting
the Stage~3 $\Delta R^2$ from marginal to significant. However, this
hypothesis must be tested empirically rather than assumed.


% =============================================================================
% Bibliography
% =============================================================================
\bibliographystyle{unsrt}
\bibliography{references}

\end{document}
